
\documentclass{article}
\usepackage{colortbl}
\usepackage{makecell}
\usepackage{multirow}
\usepackage{supertabular}

\begin{document}

\newcounter{utterance}

\centering \large Interaction Transcript for game `clean\_up', experiment `2\_hard\_5obj\_ar', episode 0 with Apertus{-}v1.5{-}8B.
\vspace{24pt}

{ \footnotesize  \setcounter{utterance}{1}
\setlength{\tabcolsep}{0pt}
\begin{supertabular}{c@{$\;$}|p{.15\linewidth}@{}p{.15\linewidth}p{.15\linewidth}p{.15\linewidth}p{.15\linewidth}p{.15\linewidth}}
    \# & $\;$A & \multicolumn{4}{c}{Game Master} & $\;\:$B\\
    \hline

    \theutterance \stepcounter{utterance}  
    & & \multicolumn{4}{p{0.6\linewidth}}{
        \cellcolor[rgb]{0.9,0.9,0.9}{
            \makecell[{{p{\linewidth}}}]{
                \texttt{\tiny{[P1$\langle$GM]}}
                \texttt{أنا مدير لعبتك، وأنت تلعب لعبة تعاونية مع الشبكة التالية كلوحة لعب:} \\
\\ 
\texttt{```} \\
\texttt{ 1234567} \\
\texttt{╔══╤════╗} \\
\texttt{║◌◌│◌◌◌H║ 1} \\
\texttt{╟──┤◌◌┌─╢ 2} \\
\texttt{║CI├──┤◌║ 3} \\
\texttt{╟──┘◌◌├─╢ 4} \\
\texttt{║◌◌◌◌◌│◌║ 5} \\
\texttt{╟─┐W┌─┼─╢ 6} \\
\texttt{║◌│T│◌│◌║ 7} \\
\texttt{╚═╧═╧═╧═╝} \\
\\ 
\texttt{```} \\
\\ 
\texttt{* تعرض الحافة العلوية إحداثيات x التي تتزايد نحو اليمين، والحافة اليمنى إحداثيات y التي تتزايد نحو الأسفل.} \\
\texttt{* الكائنات التالية موضوعة عشوائياً على شبكتك: 'W', 'I', 'T', 'C', 'H'.} \\
\\ 
\\ 
\texttt{يرى اللاعب الآخر نسخة مختلفة من لوحة اللعب، حيث توضع الكائنات في مواقع عشوائية مختلفة. لا يمكنك رؤية لوحة اللاعب الآخر، ولا يمكنه رؤية لوحتك.} \\
\\ 
\\ 
\texttt{**الهدف:**} \\
\\ 
\texttt{يجب على كلا اللاعبين تحريك الكائنات على خلفية كل منهما بحيث تنتهي الكائنات المتطابقة عند الإحداثيات نفسها. عليك التواصل مع اللاعب الآخر للاتفاق على إعداد هدف مشترك.} \\
\\ 
\\ 
\texttt{**القواعد:**} \\
\\ 
\texttt{* في كل دور، يمكنك إرسال أمر واحد فقط من الأمرين التاليين:} \\
\texttt{ 1. `قل: <رسالة>`: لإرسال رسالة (كل شيء حتى فاصل السطر التالي) إلى اللاعب الآخر. سأقوم بتمريرها إلى شريكك.} \\
\texttt{2. `حرّك: <كائن>, (<X>, <Y>)`: لتحريك كائن إلى موضع جديد، حيث `<X>` هو العمود و `<Y>` هو الصف. سأخبرك بما إذا كانت حركتك ناجحة أم لا.} \\
\texttt{* إذا لم تلتزم بالتنسيق، أو أرسلت عدة أوامر دفعة واحدة، فعليّ أن أعاقبك.} \\
\texttt{* إذا تجاوز كلا اللاعبين أكثر من 12 عقوبة، فستخسران اللعبة معاً.} \\
\texttt{* من الضروري أن تتواصل مع اللاعب الآخر بشأن حالة هدفك! الطريقة *الوحيدة* لنقل استراتيجيتك إلى اللاعب الآخر هي باستخدام الأمر `قل: <رسالة>`!} \\
\\ 
\\ 
\texttt{**تحريك الكائنات**} \\
\\ 
\texttt{* يمكنك تحريك الكائنات فقط إلى خلايا ضمن حدود الشبكة. يجب أن تكون الخلية المستهدفة فارغة، أي يجب أن تحتوي فقط على الرمز '◌'.} \\
\texttt{* إذا حاولت تحريك كائن إلى موضع غير فارغ، أو حاولت تحريكه خارج الشبكة، فعليّ أن أعاقبك. ستحصل على محاولة أخرى.} \\
\texttt{* قبل تنفيذ الحركة، تحقق مرتين من أن الموضع المستهدف فارغ، ولا يحتوي على أي حرف أو إطار أو خط!} \\
\\ 
\\ 
\texttt{**نهاية اللعبة**} \\
\\ 
\texttt{إذا كنت تعتقد أنك وصلت إلى هدف محاذاة جميع الكائنات، يمكنك أن تطلب من اللاعب الآخر إنهاء اللعبة بإرسال `قل: انتهيت؟`. إذا طلب منك اللاعب الآخر إنهاء اللعبة، ورددت بـ `قل: انتهيت!`، ستنتهي اللعبة.} \\
\\ 
\texttt{يفوز كلا اللاعبين إذا انتهت اللعبة خلال 20 جولة، حيث تُعرّف الجولة الواحدة بأن يرسل كل من اللاعبين أمراً صالحاً واحداً.} \\
\\ 
\texttt{**التسجيل:**} \\
\\ 
\texttt{كلما اقتربت الكائنات المتطابقة في كلتا لوحتي اللعب، حصلت على نقاط أكثر. العقوبات تقلل نقاطك. هل يمكنك تحطيم الرقم القياسي؟} \\
\\ 
\\ 
\texttt{يرجى إرسال رسالة إلى اللاعب الآخر لبدء اللعبة!} \\
            }
        }
    }
    & \\ \\

    \theutterance \stepcounter{utterance}  
    & \multicolumn{4}{p{0.6\linewidth}}{
        \cellcolor[rgb]{0.9,0.9,0.9}{
            \makecell[{{p{\linewidth}}}]{
                \texttt{\tiny{[P1$\rangle$GM]}}
                \texttt{قل: مرحبا، بدأنا اللعبة. سنعمل على محاذاة الكائنات. سنبدأ بفهم مواقع الكائنات على لوحتي لعبنا. هل يمكنك إرسال مواقع الكائنات على لوحتي؟} \\
            }
        }
    }
    & & \\ \\

    \theutterance \stepcounter{utterance}  
    & & & \multicolumn{2}{p{0.3\linewidth}}{
        \cellcolor[rgb]{0.9,0.9,0.9}{
            \makecell[{{p{\linewidth}}}]{
                \texttt{\tiny{[GM$|$GM]}}
                \texttt{قل: مرحبا، بدأنا اللعبة. سنعمل على محاذاة الكائنات. سنبدأ بفهم مواقع الكائنات على لوحتي لعبنا. هل يمكنك إرسال مواقع الكائنات على لوحتي؟} \\
            }
        }
    }
    & & \\ \\

    \theutterance \stepcounter{utterance}  
    & & \multicolumn{4}{p{0.6\linewidth}}{
        \cellcolor[rgb]{0.9,0.9,0.9}{
            \makecell[{{p{\linewidth}}}]{
                \texttt{\tiny{[GM$\rangle$P2]}}
                \texttt{أنا مدير لعبتك، وأنت تلعب لعبة تعاونية مع الشبكة التالية كلوحة لعب:} \\
\\ 
\texttt{```} \\
\texttt{ 1234567} \\
\texttt{╔══╤════╗} \\
\texttt{║◌W│◌◌◌◌║ 1} \\
\texttt{╟──┤◌H┌─╢ 2} \\
\texttt{║◌◌├──┤◌║ 3} \\
\texttt{╟──┘◌◌├─╢ 4} \\
\texttt{║◌◌◌◌◌│◌║ 5} \\
\texttt{╟─┐C┌─┼─╢ 6} \\
\texttt{║I│◌│T│◌║ 7} \\
\texttt{╚═╧═╧═╧═╝} \\
\\ 
\texttt{```} \\
\\ 
\texttt{* تعرض الحافة العلوية إحداثيات x التي تتزايد نحو اليمين، والحافة اليمنى إحداثيات y التي تتزايد نحو الأسفل.} \\
\texttt{* الكائنات التالية موضوعة عشوائياً على شبكتك: 'W', 'I', 'T', 'C', 'H'.} \\
\\ 
\\ 
\texttt{يرى اللاعب الآخر نسخة مختلفة من لوحة اللعب، حيث توضع الكائنات في مواقع عشوائية مختلفة. لا يمكنك رؤية لوحة اللاعب الآخر، ولا يمكنه رؤية لوحتك.} \\
\\ 
\\ 
\texttt{**الهدف:**} \\
\\ 
\texttt{يجب على كلا اللاعبين تحريك الكائنات على خلفية كل منهما بحيث تنتهي الكائنات المتطابقة عند الإحداثيات نفسها. عليك التواصل مع اللاعب الآخر للاتفاق على إعداد هدف مشترك.} \\
\\ 
\\ 
\texttt{**القواعد:**} \\
\\ 
\texttt{* في كل دور، يمكنك إرسال أمر واحد فقط من الأمرين التاليين:} \\
\texttt{ 1. `قل: <رسالة>`: لإرسال رسالة (كل شيء حتى فاصل السطر التالي) إلى اللاعب الآخر. سأقوم بتمريرها إلى شريكك.} \\
\texttt{2. `حرّك: <كائن>, (<X>, <Y>)`: لتحريك كائن إلى موضع جديد، حيث `<X>` هو العمود و `<Y>` هو الصف. سأخبرك بما إذا كانت حركتك ناجحة أم لا.} \\
\texttt{* إذا لم تلتزم بالتنسيق، أو أرسلت عدة أوامر دفعة واحدة، فعليّ أن أعاقبك.} \\
\texttt{* إذا تجاوز كلا اللاعبين أكثر من 12 عقوبة، فستخسران اللعبة معاً.} \\
\texttt{* من الضروري أن تتواصل مع اللاعب الآخر بشأن حالة هدفك! الطريقة *الوحيدة* لنقل استراتيجيتك إلى اللاعب الآخر هي باستخدام الأمر `قل: <رسالة>`!} \\
\\ 
\\ 
\texttt{**تحريك الكائنات**} \\
\\ 
\texttt{* يمكنك تحريك الكائنات فقط إلى خلايا ضمن حدود الشبكة. يجب أن تكون الخلية المستهدفة فارغة، أي يجب أن تحتوي فقط على الرمز '◌'.} \\
\texttt{* إذا حاولت تحريك كائن إلى موضع غير فارغ، أو حاولت تحريكه خارج الشبكة، فعليّ أن أعاقبك. ستحصل على محاولة أخرى.} \\
\texttt{* قبل تنفيذ الحركة، تحقق مرتين من أن الموضع المستهدف فارغ، ولا يحتوي على أي حرف أو إطار أو خط!} \\
\\ 
\\ 
\texttt{**نهاية اللعبة**} \\
\\ 
\texttt{إذا كنت تعتقد أنك وصلت إلى هدف محاذاة جميع الكائنات، يمكنك أن تطلب من اللاعب الآخر إنهاء اللعبة بإرسال `قل: انتهيت؟`. إذا طلب منك اللاعب الآخر إنهاء اللعبة، ورددت بـ `قل: انتهيت!`، ستنتهي اللعبة.} \\
\\ 
\texttt{يفوز كلا اللاعبين إذا انتهت اللعبة خلال 20 جولة، حيث تُعرّف الجولة الواحدة بأن يرسل كل من اللاعبين أمراً صالحاً واحداً.} \\
\\ 
\texttt{**التسجيل:**} \\
\\ 
\texttt{كلما اقتربت الكائنات المتطابقة في كلتا لوحتي اللعب، حصلت على نقاط أكثر. العقوبات تقلل نقاطك. هل يمكنك تحطيم الرقم القياسي؟} \\
\\ 
\texttt{بدأ اللاعب الآخر اللعبة بإرسال هذه الرسالة:} \\
\\ 
\texttt{"مرحبا، بدأنا اللعبة. سنعمل على محاذاة الكائنات. سنبدأ بفهم مواقع الكائنات على لوحتي لعبنا. هل يمكنك إرسال مواقع الكائنات على لوحتي؟"} \\
\\ 
\texttt{ما هو أمرك الأول؟} \\
            }
        }
    }
    & \\ \\

    \theutterance \stepcounter{utterance}  
    & & & \multicolumn{4}{p{0.6\linewidth}}{
        \cellcolor[rgb]{0.9,0.9,0.9}{
            \makecell[{{p{\linewidth}}}]{
                \texttt{\tiny{[GM$\langle$P2]}}
                \texttt{قل: سأرسل مواقع الكائنات على لوحتي. سأقوم بتوثيق كل كائن مع إحداثياته (X,Y) بناءً على شبكة العرض، وأرسلها لك.} \\
            }
        }
    }
     \\ \\

    \theutterance \stepcounter{utterance}  
    & & & \multicolumn{2}{p{0.3\linewidth}}{
        \cellcolor[rgb]{0.9,0.9,0.9}{
            \makecell[{{p{\linewidth}}}]{
                \texttt{\tiny{[GM$|$GM]}}
                \texttt{قل: سأرسل مواقع الكائنات على لوحتي. سأقوم بتوثيق كل كائن مع إحداثياته (X,Y) بناءً على شبكة العرض، وأرسلها لك.} \\
            }
        }
    }
    & & \\ \\

    \theutterance \stepcounter{utterance}  
    & & \multicolumn{4}{p{0.6\linewidth}}{
        \cellcolor[rgb]{0.9,0.9,0.9}{
            \makecell[{{p{\linewidth}}}]{
                \texttt{\tiny{[P1$\langle$GM]}}
                \texttt{تم تمرير رسالتك إلى اللاعب الآخر.} \\
\\ 
\\ 
\texttt{أنت تلعب حالياً الجولة 1 من أصل 20 كحد أقصى.} \\
\texttt{لقد جمعتما معاً 0 من أصل 12 عقوبة.} \\
\texttt{أرسل اللاعب الآخر هذه الرسالة:} \\
\\ 
\\ 
\texttt{"سأرسل مواقع الكائنات على لوحتي. سأقوم بتوثيق كل كائن مع إحداثياته (X,Y) بناءً على شبكة العرض، وأرسلها لك."} \\
\\ 
\\ 
\texttt{ما هو أمرك التالي؟} \\
            }
        }
    }
    & \\ \\

    \theutterance \stepcounter{utterance}  
    & \multicolumn{4}{p{0.6\linewidth}}{
        \cellcolor[rgb]{0.9,0.9,0.9}{
            \makecell[{{p{\linewidth}}}]{
                \texttt{\tiny{[P1$\rangle$GM]}}
                \texttt{قل: حسناً، سأعدّ جدولاً بالأحداث. يرجى إرسال مواقع الكائنات مع إحداثياتها (X,Y) بناءً على الشبكة، وسأقوم بردم هذه المعلومات لتخطيط الخطوات التالية.} \\
            }
        }
    }
    & & \\ \\

    \theutterance \stepcounter{utterance}  
    & & & \multicolumn{2}{p{0.3\linewidth}}{
        \cellcolor[rgb]{0.9,0.9,0.9}{
            \makecell[{{p{\linewidth}}}]{
                \texttt{\tiny{[GM$|$GM]}}
                \texttt{قل: حسناً، سأعدّ جدولاً بالأحداث. يرجى إرسال مواقع الكائنات مع إحداثياتها (X,Y) بناءً على الشبكة، وسأقوم بردم هذه المعلومات لتخطيط الخطوات التالية.} \\
            }
        }
    }
    & & \\ \\

    \theutterance \stepcounter{utterance}  
    & & \multicolumn{4}{p{0.6\linewidth}}{
        \cellcolor[rgb]{0.9,0.9,0.9}{
            \makecell[{{p{\linewidth}}}]{
                \texttt{\tiny{[GM$\rangle$P2]}}
                \texttt{تم تمرير رسالتك إلى اللاعب الآخر.} \\
\\ 
\\ 
\texttt{أنت تلعب حالياً الجولة 2 من أصل 20 كحد أقصى.} \\
\texttt{لقد جمعتما معاً 0 من أصل 12 عقوبة.} \\
\texttt{أرسل اللاعب الآخر هذه الرسالة:} \\
\\ 
\\ 
\texttt{"حسناً، سأعدّ جدولاً بالأحداث. يرجى إرسال مواقع الكائنات مع إحداثياتها (X,Y) بناءً على الشبكة، وسأقوم بردم هذه المعلومات لتخطيط الخطوات التالية."} \\
\\ 
\\ 
\texttt{ما هو أمرك التالي؟} \\
            }
        }
    }
    & \\ \\

    \theutterance \stepcounter{utterance}  
    & & & \multicolumn{4}{p{0.6\linewidth}}{
        \cellcolor[rgb]{0.9,0.9,0.9}{
            \makecell[{{p{\linewidth}}}]{
                \texttt{\tiny{[GM$\langle$P2]}}
                \texttt{قل: سأقوم بتوثيق مواقع الكائنات على لوحتي وسأرسل إحداثياتها مع توضيح كل كائن. سأبدأ بتحديد كل كائن مع X و Y كما تظهر على لوحتي، ثم أرسلها لك.} \\
            }
        }
    }
     \\ \\

    \theutterance \stepcounter{utterance}  
    & & & \multicolumn{2}{p{0.3\linewidth}}{
        \cellcolor[rgb]{0.9,0.9,0.9}{
            \makecell[{{p{\linewidth}}}]{
                \texttt{\tiny{[GM$|$GM]}}
                \texttt{قل: سأقوم بتوثيق مواقع الكائنات على لوحتي وسأرسل إحداثياتها مع توضيح كل كائن. سأبدأ بتحديد كل كائن مع X و Y كما تظهر على لوحتي، ثم أرسلها لك.} \\
            }
        }
    }
    & & \\ \\

    \theutterance \stepcounter{utterance}  
    & & \multicolumn{4}{p{0.6\linewidth}}{
        \cellcolor[rgb]{0.9,0.9,0.9}{
            \makecell[{{p{\linewidth}}}]{
                \texttt{\tiny{[P1$\langle$GM]}}
                \texttt{تم تمرير رسالتك إلى اللاعب الآخر.} \\
\\ 
\\ 
\texttt{أنت تلعب حالياً الجولة 2 من أصل 20 كحد أقصى.} \\
\texttt{لقد جمعتما معاً 0 من أصل 12 عقوبة.} \\
\texttt{أرسل اللاعب الآخر هذه الرسالة:} \\
\\ 
\\ 
\texttt{"سأقوم بتوثيق مواقع الكائنات على لوحتي وسأرسل إحداثياتها مع توضيح كل كائن. سأبدأ بتحديد كل كائن مع X و Y كما تظهر على لوحتي، ثم أرسلها لك."} \\
\\ 
\\ 
\texttt{ما هو أمرك التالي؟} \\
            }
        }
    }
    & \\ \\

    \theutterance \stepcounter{utterance}  
    & \multicolumn{4}{p{0.6\linewidth}}{
        \cellcolor[rgb]{0.9,0.9,0.9}{
            \makecell[{{p{\linewidth}}}]{
                \texttt{\tiny{[P1$\rangle$GM]}}
                \texttt{قل: حسناً، انتظر الرسالة التي تحتوي على توثيق مواقع الكائنات مع إحداثيات X و Y لكل كائن. سأقوم باستقبال المعلومات ومعالجتها فوراً، ثم سأحدد الخطوات التالية في المحاذاة.} \\
            }
        }
    }
    & & \\ \\

    \theutterance \stepcounter{utterance}  
    & & & \multicolumn{2}{p{0.3\linewidth}}{
        \cellcolor[rgb]{0.9,0.9,0.9}{
            \makecell[{{p{\linewidth}}}]{
                \texttt{\tiny{[GM$|$GM]}}
                \texttt{قل: حسناً، انتظر الرسالة التي تحتوي على توثيق مواقع الكائنات مع إحداثيات X و Y لكل كائن. سأقوم باستقبال المعلومات ومعالجتها فوراً، ثم سأحدد الخطوات التالية في المحاذاة.} \\
            }
        }
    }
    & & \\ \\

    \theutterance \stepcounter{utterance}  
    & & \multicolumn{4}{p{0.6\linewidth}}{
        \cellcolor[rgb]{0.9,0.9,0.9}{
            \makecell[{{p{\linewidth}}}]{
                \texttt{\tiny{[GM$\rangle$P2]}}
                \texttt{تم تمرير رسالتك إلى اللاعب الآخر.} \\
\\ 
\\ 
\texttt{أنت تلعب حالياً الجولة 3 من أصل 20 كحد أقصى.} \\
\texttt{لقد جمعتما معاً 0 من أصل 12 عقوبة.} \\
\texttt{أرسل اللاعب الآخر هذه الرسالة:} \\
\\ 
\\ 
\texttt{"حسناً، انتظر الرسالة التي تحتوي على توثيق مواقع الكائنات مع إحداثيات X و Y لكل كائن. سأقوم باستقبال المعلومات ومعالجتها فوراً، ثم سأحدد الخطوات التالية في المحاذاة."} \\
\\ 
\\ 
\texttt{ما هو أمرك التالي؟} \\
            }
        }
    }
    & \\ \\

    \theutterance \stepcounter{utterance}  
    & & & \multicolumn{4}{p{0.6\linewidth}}{
        \cellcolor[rgb]{0.9,0.9,0.9}{
            \makecell[{{p{\linewidth}}}]{
                \texttt{\tiny{[GM$\langle$P2]}}
                \texttt{حرّك: (سيتم إرسال موقع كائن..)} \\
            }
        }
    }
     \\ \\

    \theutterance \stepcounter{utterance}  
    & & & \multicolumn{2}{p{0.3\linewidth}}{
        \cellcolor[rgb]{0.9,0.9,0.9}{
            \makecell[{{p{\linewidth}}}]{
                \texttt{\tiny{[GM$|$GM]}}
                \texttt{حرّك: (سيتم إرسال موقع كائن..)} \\
            }
        }
    }
    & & \\ \\

    \theutterance \stepcounter{utterance}  
    & & & \multicolumn{2}{p{0.3\linewidth}}{
        \cellcolor[rgb]{0.9,0.9,0.9}{
            \makecell[{{p{\linewidth}}}]{
                \texttt{\tiny{[GM$|$GM]}}
                \texttt{Invalid response format} \\
            }
        }
    }
    & & \\ \\

    \theutterance \stepcounter{utterance}  
    & & \multicolumn{4}{p{0.6\linewidth}}{
        \cellcolor[rgb]{0.9,0.9,0.9}{
            \makecell[{{p{\linewidth}}}]{
                \texttt{\tiny{[GM$\rangle$P2]}}
                \texttt{عقوبة: رسالتك ليست بالتنسيق المتوقع!} \\
\\ 
\texttt{تأكد من أن ردك يحتوي فقط على `قل: <رسالة>` أو `حرّك: <كائن>, (<X>, <Y>)`، ولا شيء غير ذلك!} \\
\\ 
\texttt{لقد جمعتما معاً 1 من أصل 12 عقوبة.} \\
\texttt{يرجى المحاولة مرة أخرى!} \\
            }
        }
    }
    & \\ \\

    \theutterance \stepcounter{utterance}  
    & & & \multicolumn{4}{p{0.6\linewidth}}{
        \cellcolor[rgb]{0.9,0.9,0.9}{
            \makecell[{{p{\linewidth}}}]{
                \texttt{\tiny{[GM$\langle$P2]}}
                \texttt{قل: سأقوم بإرسال توثيق مواقع الكائنات مع إحداثيات X و Y لكل منها، وفقاً للشبكة المعروضة حالياً على لوحتي.} \\
            }
        }
    }
     \\ \\

    \theutterance \stepcounter{utterance}  
    & & & \multicolumn{2}{p{0.3\linewidth}}{
        \cellcolor[rgb]{0.9,0.9,0.9}{
            \makecell[{{p{\linewidth}}}]{
                \texttt{\tiny{[GM$|$GM]}}
                \texttt{قل: سأقوم بإرسال توثيق مواقع الكائنات مع إحداثيات X و Y لكل منها، وفقاً للشبكة المعروضة حالياً على لوحتي.} \\
            }
        }
    }
    & & \\ \\

    \theutterance \stepcounter{utterance}  
    & & \multicolumn{4}{p{0.6\linewidth}}{
        \cellcolor[rgb]{0.9,0.9,0.9}{
            \makecell[{{p{\linewidth}}}]{
                \texttt{\tiny{[P1$\langle$GM]}}
                \texttt{تم تمرير رسالتك إلى اللاعب الآخر.} \\
\\ 
\\ 
\texttt{أنت تلعب حالياً الجولة 3 من أصل 20 كحد أقصى.} \\
\texttt{لقد جمعتما معاً 1 من أصل 12 عقوبة.} \\
\texttt{أرسل اللاعب الآخر هذه الرسالة:} \\
\\ 
\\ 
\texttt{"سأقوم بإرسال توثيق مواقع الكائنات مع إحداثيات X و Y لكل منها، وفقاً للشبكة المعروضة حالياً على لوحتي."} \\
\\ 
\\ 
\texttt{ما هو أمرك التالي؟} \\
            }
        }
    }
    & \\ \\

    \theutterance \stepcounter{utterance}  
    & \multicolumn{4}{p{0.6\linewidth}}{
        \cellcolor[rgb]{0.9,0.9,0.9}{
            \makecell[{{p{\linewidth}}}]{
                \texttt{\tiny{[P1$\rangle$GM]}}
                \texttt{قل: يرجى إرسال توثيق مواقع الكائنات مع إحداثيات X و Y لكل منها، بناءً على الشبكة المعروضة حالياً على لوحتي. سأقوم باستقبال المعلومات ومعالجتها لاستكشاف الخطة التالية.} \\
            }
        }
    }
    & & \\ \\

    \theutterance \stepcounter{utterance}  
    & & & \multicolumn{2}{p{0.3\linewidth}}{
        \cellcolor[rgb]{0.9,0.9,0.9}{
            \makecell[{{p{\linewidth}}}]{
                \texttt{\tiny{[GM$|$GM]}}
                \texttt{قل: يرجى إرسال توثيق مواقع الكائنات مع إحداثيات X و Y لكل منها، بناءً على الشبكة المعروضة حالياً على لوحتي. سأقوم باستقبال المعلومات ومعالجتها لاستكشاف الخطة التالية.} \\
            }
        }
    }
    & & \\ \\

    \theutterance \stepcounter{utterance}  
    & & \multicolumn{4}{p{0.6\linewidth}}{
        \cellcolor[rgb]{0.9,0.9,0.9}{
            \makecell[{{p{\linewidth}}}]{
                \texttt{\tiny{[GM$\rangle$P2]}}
                \texttt{تم تمرير رسالتك إلى اللاعب الآخر.} \\
\\ 
\\ 
\texttt{أنت تلعب حالياً الجولة 4 من أصل 20 كحد أقصى.} \\
\texttt{لقد جمعتما معاً 1 من أصل 12 عقوبة.} \\
\texttt{أرسل اللاعب الآخر هذه الرسالة:} \\
\\ 
\\ 
\texttt{"يرجى إرسال توثيق مواقع الكائنات مع إحداثيات X و Y لكل منها، بناءً على الشبكة المعروضة حالياً على لوحتي. سأقوم باستقبال المعلومات ومعالجتها لاستكشاف الخطة التالية."} \\
\\ 
\\ 
\texttt{ما هو أمرك التالي؟} \\
            }
        }
    }
    & \\ \\

    \theutterance \stepcounter{utterance}  
    & & & \multicolumn{4}{p{0.6\linewidth}}{
        \cellcolor[rgb]{0.9,0.9,0.9}{
            \makecell[{{p{\linewidth}}}]{
                \texttt{\tiny{[GM$\langle$P2]}}
                \texttt{حرّك: أنا أشرح حالياً توثيق المواقع، لذا سأستخدم الأمر 'قل' للتقديم: قل: سأقوم بتوثيق كل كائن بوضعه وإحداثياته (X,Y) كما تظهر على لوحتي حالياً، ثم سأرسلها لك في الخطوة التالية.} \\
            }
        }
    }
     \\ \\

    \theutterance \stepcounter{utterance}  
    & & & \multicolumn{2}{p{0.3\linewidth}}{
        \cellcolor[rgb]{0.9,0.9,0.9}{
            \makecell[{{p{\linewidth}}}]{
                \texttt{\tiny{[GM$|$GM]}}
                \texttt{حرّك: أنا أشرح حالياً توثيق المواقع، لذا سأستخدم الأمر 'قل' للتقديم: قل: سأقوم بتوثيق كل كائن بوضعه وإحداثياته (X,Y) كما تظهر على لوحتي حالياً، ثم سأرسلها لك في الخطوة التالية.} \\
            }
        }
    }
    & & \\ \\

    \theutterance \stepcounter{utterance}  
    & & & \multicolumn{2}{p{0.3\linewidth}}{
        \cellcolor[rgb]{0.9,0.9,0.9}{
            \makecell[{{p{\linewidth}}}]{
                \texttt{\tiny{[GM$|$GM]}}
                \texttt{Invalid format: head is not empty: 'حرّك: أنا أشرح حالياً توثيق المواقع، لذا سأستخدم الأمر 'قل' للتقديم: '} \\
            }
        }
    }
    & & \\ \\

    \theutterance \stepcounter{utterance}  
    & & \multicolumn{4}{p{0.6\linewidth}}{
        \cellcolor[rgb]{0.9,0.9,0.9}{
            \makecell[{{p{\linewidth}}}]{
                \texttt{\tiny{[GM$\rangle$P2]}}
                \texttt{عقوبة: يجب ألا تحتوي رسالتك على أي شيء قبل الأمر!} \\
\\ 
\texttt{تأكد من أن ردك يحتوي فقط على `قل: <رسالة>` أو `حرّك: <كائن>, (<X>, <Y>)`، ولا شيء غير ذلك!} \\
\\ 
\texttt{لقد جمعتما معاً 2 من أصل 12 عقوبة.} \\
\texttt{يرجى المحاولة مرة أخرى!} \\
            }
        }
    }
    & \\ \\

    \theutterance \stepcounter{utterance}  
    & & & \multicolumn{4}{p{0.6\linewidth}}{
        \cellcolor[rgb]{0.9,0.9,0.9}{
            \makecell[{{p{\linewidth}}}]{
                \texttt{\tiny{[GM$\langle$P2]}}
                \texttt{قل: سأقوم بتوثيق كل كائن بوضعه وإحداثياته (X,Y) كما تظهر على لوحتي حالياً، وسأرسلها لك خلال الدورية التالية.} \\
            }
        }
    }
     \\ \\

    \theutterance \stepcounter{utterance}  
    & & & \multicolumn{2}{p{0.3\linewidth}}{
        \cellcolor[rgb]{0.9,0.9,0.9}{
            \makecell[{{p{\linewidth}}}]{
                \texttt{\tiny{[GM$|$GM]}}
                \texttt{قل: سأقوم بتوثيق كل كائن بوضعه وإحداثياته (X,Y) كما تظهر على لوحتي حالياً، وسأرسلها لك خلال الدورية التالية.} \\
            }
        }
    }
    & & \\ \\

    \theutterance \stepcounter{utterance}  
    & & \multicolumn{4}{p{0.6\linewidth}}{
        \cellcolor[rgb]{0.9,0.9,0.9}{
            \makecell[{{p{\linewidth}}}]{
                \texttt{\tiny{[P1$\langle$GM]}}
                \texttt{تم تمرير رسالتك إلى اللاعب الآخر.} \\
\\ 
\\ 
\texttt{أنت تلعب حالياً الجولة 4 من أصل 20 كحد أقصى.} \\
\texttt{لقد جمعتما معاً 2 من أصل 12 عقوبة.} \\
\texttt{أرسل اللاعب الآخر هذه الرسالة:} \\
\\ 
\\ 
\texttt{"سأقوم بتوثيق كل كائن بوضعه وإحداثياته (X,Y) كما تظهر على لوحتي حالياً، وسأرسلها لك خلال الدورية التالية."} \\
\\ 
\\ 
\texttt{ما هو أمرك التالي؟} \\
            }
        }
    }
    & \\ \\

    \theutterance \stepcounter{utterance}  
    & \multicolumn{4}{p{0.6\linewidth}}{
        \cellcolor[rgb]{0.9,0.9,0.9}{
            \makecell[{{p{\linewidth}}}]{
                \texttt{\tiny{[P1$\rangle$GM]}}
                \texttt{قل: حسناً، أشكرك على التوضيح. سأعدّ لأستقبال قائمة توثيق المواقع مع إحداثيات X و Y لكل كائن في الدورية القادمة. سأتوقع الرسالة خلال الدورية التالية، وأكون جاهزاً لتحليل المعلومات لتعريف الخطة المشتركة للمحاذاة.} \\
            }
        }
    }
    & & \\ \\

    \theutterance \stepcounter{utterance}  
    & & & \multicolumn{2}{p{0.3\linewidth}}{
        \cellcolor[rgb]{0.9,0.9,0.9}{
            \makecell[{{p{\linewidth}}}]{
                \texttt{\tiny{[GM$|$GM]}}
                \texttt{قل: حسناً، أشكرك على التوضيح. سأعدّ لأستقبال قائمة توثيق المواقع مع إحداثيات X و Y لكل كائن في الدورية القادمة. سأتوقع الرسالة خلال الدورية التالية، وأكون جاهزاً لتحليل المعلومات لتعريف الخطة المشتركة للمحاذاة.} \\
            }
        }
    }
    & & \\ \\

    \theutterance \stepcounter{utterance}  
    & & \multicolumn{4}{p{0.6\linewidth}}{
        \cellcolor[rgb]{0.9,0.9,0.9}{
            \makecell[{{p{\linewidth}}}]{
                \texttt{\tiny{[GM$\rangle$P2]}}
                \texttt{تم تمرير رسالتك إلى اللاعب الآخر.} \\
\\ 
\\ 
\texttt{أنت تلعب حالياً الجولة 5 من أصل 20 كحد أقصى.} \\
\texttt{لقد جمعتما معاً 2 من أصل 12 عقوبة.} \\
\texttt{أرسل اللاعب الآخر هذه الرسالة:} \\
\\ 
\\ 
\texttt{"حسناً، أشكرك على التوضيح. سأعدّ لأستقبال قائمة توثيق المواقع مع إحداثيات X و Y لكل كائن في الدورية القادمة. سأتوقع الرسالة خلال الدورية التالية، وأكون جاهزاً لتحليل المعلومات لتعريف الخطة المشتركة للمحاذاة."} \\
\\ 
\\ 
\texttt{ما هو أمرك التالي؟} \\
            }
        }
    }
    & \\ \\

    \theutterance \stepcounter{utterance}  
    & & & \multicolumn{4}{p{0.6\linewidth}}{
        \cellcolor[rgb]{0.9,0.9,0.9}{
            \makecell[{{p{\linewidth}}}]{
                \texttt{\tiny{[GM$\langle$P2]}}
                \texttt{قل: سأقوم بإرسال توثيق دقيق ومفصل لكل كائن، حيث سأذكر اسم الكائن مع إحداثياته (X,Y) كما تظهر على لوحتي حالياً، ثم سأتنازل لك ليقوم بتحليلها.} \\
            }
        }
    }
     \\ \\

    \theutterance \stepcounter{utterance}  
    & & & \multicolumn{2}{p{0.3\linewidth}}{
        \cellcolor[rgb]{0.9,0.9,0.9}{
            \makecell[{{p{\linewidth}}}]{
                \texttt{\tiny{[GM$|$GM]}}
                \texttt{قل: سأقوم بإرسال توثيق دقيق ومفصل لكل كائن، حيث سأذكر اسم الكائن مع إحداثياته (X,Y) كما تظهر على لوحتي حالياً، ثم سأتنازل لك ليقوم بتحليلها.} \\
            }
        }
    }
    & & \\ \\

    \theutterance \stepcounter{utterance}  
    & & \multicolumn{4}{p{0.6\linewidth}}{
        \cellcolor[rgb]{0.9,0.9,0.9}{
            \makecell[{{p{\linewidth}}}]{
                \texttt{\tiny{[P1$\langle$GM]}}
                \texttt{تم تمرير رسالتك إلى اللاعب الآخر.} \\
\\ 
\\ 
\texttt{أنت تلعب حالياً الجولة 5 من أصل 20 كحد أقصى.} \\
\texttt{لقد جمعتما معاً 2 من أصل 12 عقوبة.} \\
\texttt{أرسل اللاعب الآخر هذه الرسالة:} \\
\\ 
\\ 
\texttt{"سأقوم بإرسال توثيق دقيق ومفصل لكل كائن، حيث سأذكر اسم الكائن مع إحداثياته (X,Y) كما تظهر على لوحتي حالياً، ثم سأتنازل لك ليقوم بتحليلها."} \\
\\ 
\\ 
\texttt{ما هو أمرك التالي؟} \\
            }
        }
    }
    & \\ \\

    \theutterance \stepcounter{utterance}  
    & \multicolumn{4}{p{0.6\linewidth}}{
        \cellcolor[rgb]{0.9,0.9,0.9}{
            \makecell[{{p{\linewidth}}}]{
                \texttt{\tiny{[P1$\rangle$GM]}}
                \texttt{قل: شكراً لك على الرسالة. انتظر توثيق كل كائن مع اسمه وأحداثياته (X,Y) كما تظهر على لوحتي. سأتولى تحليل هذه المعلومات فوراً لتحديد المواقع الصحيحة لكل كائن ثم نحدد الخطوات المشتركة للمحاذاة.} \\
            }
        }
    }
    & & \\ \\

    \theutterance \stepcounter{utterance}  
    & & & \multicolumn{2}{p{0.3\linewidth}}{
        \cellcolor[rgb]{0.9,0.9,0.9}{
            \makecell[{{p{\linewidth}}}]{
                \texttt{\tiny{[GM$|$GM]}}
                \texttt{قل: شكراً لك على الرسالة. انتظر توثيق كل كائن مع اسمه وأحداثياته (X,Y) كما تظهر على لوحتي. سأتولى تحليل هذه المعلومات فوراً لتحديد المواقع الصحيحة لكل كائن ثم نحدد الخطوات المشتركة للمحاذاة.} \\
            }
        }
    }
    & & \\ \\

    \theutterance \stepcounter{utterance}  
    & & \multicolumn{4}{p{0.6\linewidth}}{
        \cellcolor[rgb]{0.9,0.9,0.9}{
            \makecell[{{p{\linewidth}}}]{
                \texttt{\tiny{[GM$\rangle$P2]}}
                \texttt{تم تمرير رسالتك إلى اللاعب الآخر.} \\
\\ 
\\ 
\texttt{أنت تلعب حالياً الجولة 6 من أصل 20 كحد أقصى.} \\
\texttt{لقد جمعتما معاً 2 من أصل 12 عقوبة.} \\
\texttt{أرسل اللاعب الآخر هذه الرسالة:} \\
\\ 
\\ 
\texttt{"شكراً لك على الرسالة. انتظر توثيق كل كائن مع اسمه وأحداثياته (X,Y) كما تظهر على لوحتي. سأتولى تحليل هذه المعلومات فوراً لتحديد المواقع الصحيحة لكل كائن ثم نحدد الخطوات المشتركة للمحاذاة."} \\
\\ 
\\ 
\texttt{ما هو أمرك التالي؟} \\
            }
        }
    }
    & \\ \\

    \theutterance \stepcounter{utterance}  
    & & & \multicolumn{4}{p{0.6\linewidth}}{
        \cellcolor[rgb]{0.9,0.9,0.9}{
            \makecell[{{p{\linewidth}}}]{
                \texttt{\tiny{[GM$\langle$P2]}}
                \texttt{حرّك: (سيتم إرسال توثيق الكائنات..)} \\
            }
        }
    }
     \\ \\

    \theutterance \stepcounter{utterance}  
    & & & \multicolumn{2}{p{0.3\linewidth}}{
        \cellcolor[rgb]{0.9,0.9,0.9}{
            \makecell[{{p{\linewidth}}}]{
                \texttt{\tiny{[GM$|$GM]}}
                \texttt{حرّك: (سيتم إرسال توثيق الكائنات..)} \\
            }
        }
    }
    & & \\ \\

    \theutterance \stepcounter{utterance}  
    & & & \multicolumn{2}{p{0.3\linewidth}}{
        \cellcolor[rgb]{0.9,0.9,0.9}{
            \makecell[{{p{\linewidth}}}]{
                \texttt{\tiny{[GM$|$GM]}}
                \texttt{Invalid response format} \\
            }
        }
    }
    & & \\ \\

    \theutterance \stepcounter{utterance}  
    & & \multicolumn{4}{p{0.6\linewidth}}{
        \cellcolor[rgb]{0.9,0.9,0.9}{
            \makecell[{{p{\linewidth}}}]{
                \texttt{\tiny{[GM$\rangle$P2]}}
                \texttt{عقوبة: رسالتك ليست بالتنسيق المتوقع!} \\
\\ 
\texttt{تأكد من أن ردك يحتوي فقط على `قل: <رسالة>` أو `حرّك: <كائن>, (<X>, <Y>)`، ولا شيء غير ذلك!} \\
\\ 
\texttt{لقد جمعتما معاً 3 من أصل 12 عقوبة.} \\
\texttt{يرجى المحاولة مرة أخرى!} \\
            }
        }
    }
    & \\ \\

    \theutterance \stepcounter{utterance}  
    & & & \multicolumn{4}{p{0.6\linewidth}}{
        \cellcolor[rgb]{0.9,0.9,0.9}{
            \makecell[{{p{\linewidth}}}]{
                \texttt{\tiny{[GM$\langle$P2]}}
                \texttt{قل: سأقوم بتوثيق كل كائن مع وضعه وأحداثياته (X,Y) كما تظهر على لوحتي حالياً، ثم سأرسل التفاصيل الدقيقة لك خلال الدورية التالية.} \\
            }
        }
    }
     \\ \\

    \theutterance \stepcounter{utterance}  
    & & & \multicolumn{2}{p{0.3\linewidth}}{
        \cellcolor[rgb]{0.9,0.9,0.9}{
            \makecell[{{p{\linewidth}}}]{
                \texttt{\tiny{[GM$|$GM]}}
                \texttt{قل: سأقوم بتوثيق كل كائن مع وضعه وأحداثياته (X,Y) كما تظهر على لوحتي حالياً، ثم سأرسل التفاصيل الدقيقة لك خلال الدورية التالية.} \\
            }
        }
    }
    & & \\ \\

    \theutterance \stepcounter{utterance}  
    & & \multicolumn{4}{p{0.6\linewidth}}{
        \cellcolor[rgb]{0.9,0.9,0.9}{
            \makecell[{{p{\linewidth}}}]{
                \texttt{\tiny{[P1$\langle$GM]}}
                \texttt{تم تمرير رسالتك إلى اللاعب الآخر.} \\
\\ 
\\ 
\texttt{أنت تلعب حالياً الجولة 6 من أصل 20 كحد أقصى.} \\
\texttt{لقد جمعتما معاً 3 من أصل 12 عقوبة.} \\
\texttt{أرسل اللاعب الآخر هذه الرسالة:} \\
\\ 
\\ 
\texttt{"سأقوم بتوثيق كل كائن مع وضعه وأحداثياته (X,Y) كما تظهر على لوحتي حالياً، ثم سأرسل التفاصيل الدقيقة لك خلال الدورية التالية."} \\
\\ 
\\ 
\texttt{ما هو أمرك التالي؟} \\
            }
        }
    }
    & \\ \\

    \theutterance \stepcounter{utterance}  
    & \multicolumn{4}{p{0.6\linewidth}}{
        \cellcolor[rgb]{0.9,0.9,0.9}{
            \makecell[{{p{\linewidth}}}]{
                \texttt{\tiny{[P1$\rangle$GM]}}
                \texttt{قل: حسناً، أشكرك على التوضيح. انتظر التفاصيل الدقيقة لكل كائن مع اسمه وأحداثياته (X,Y) كما تظهر على لوحتي، وسأتولى التحليل فوراً لتحديد المواقع الصحيحة لكل كائن وبدء الخطة المشتركة للمحاذاة.} \\
            }
        }
    }
    & & \\ \\

    \theutterance \stepcounter{utterance}  
    & & & \multicolumn{2}{p{0.3\linewidth}}{
        \cellcolor[rgb]{0.9,0.9,0.9}{
            \makecell[{{p{\linewidth}}}]{
                \texttt{\tiny{[GM$|$GM]}}
                \texttt{قل: حسناً، أشكرك على التوضيح. انتظر التفاصيل الدقيقة لكل كائن مع اسمه وأحداثياته (X,Y) كما تظهر على لوحتي، وسأتولى التحليل فوراً لتحديد المواقع الصحيحة لكل كائن وبدء الخطة المشتركة للمحاذاة.} \\
            }
        }
    }
    & & \\ \\

    \theutterance \stepcounter{utterance}  
    & & \multicolumn{4}{p{0.6\linewidth}}{
        \cellcolor[rgb]{0.9,0.9,0.9}{
            \makecell[{{p{\linewidth}}}]{
                \texttt{\tiny{[GM$\rangle$P2]}}
                \texttt{تم تمرير رسالتك إلى اللاعب الآخر.} \\
\\ 
\\ 
\texttt{أنت تلعب حالياً الجولة 7 من أصل 20 كحد أقصى.} \\
\texttt{لقد جمعتما معاً 3 من أصل 12 عقوبة.} \\
\texttt{أرسل اللاعب الآخر هذه الرسالة:} \\
\\ 
\\ 
\texttt{"حسناً، أشكرك على التوضيح. انتظر التفاصيل الدقيقة لكل كائن مع اسمه وأحداثياته (X,Y) كما تظهر على لوحتي، وسأتولى التحليل فوراً لتحديد المواقع الصحيحة لكل كائن وبدء الخطة المشتركة للمحاذاة."} \\
\\ 
\\ 
\texttt{ما هو أمرك التالي؟} \\
            }
        }
    }
    & \\ \\

    \theutterance \stepcounter{utterance}  
    & & & \multicolumn{4}{p{0.6\linewidth}}{
        \cellcolor[rgb]{0.9,0.9,0.9}{
            \makecell[{{p{\linewidth}}}]{
                \texttt{\tiny{[GM$\langle$P2]}}
                \texttt{قل: سأقوم بتوثيق كل كائن بوضعه وإحداثياته (X,Y) كما تظهر على لوحتي: W في (1,3), H في (2,2), I في (1,7), C في (5,6), T في (7,5).} \\
            }
        }
    }
     \\ \\

    \theutterance \stepcounter{utterance}  
    & & & \multicolumn{2}{p{0.3\linewidth}}{
        \cellcolor[rgb]{0.9,0.9,0.9}{
            \makecell[{{p{\linewidth}}}]{
                \texttt{\tiny{[GM$|$GM]}}
                \texttt{قل: سأقوم بتوثيق كل كائن بوضعه وإحداثياته (X,Y) كما تظهر على لوحتي: W في (1,3), H في (2,2), I في (1,7), C في (5,6), T في (7,5).} \\
            }
        }
    }
    & & \\ \\

    \theutterance \stepcounter{utterance}  
    & & \multicolumn{4}{p{0.6\linewidth}}{
        \cellcolor[rgb]{0.9,0.9,0.9}{
            \makecell[{{p{\linewidth}}}]{
                \texttt{\tiny{[P1$\langle$GM]}}
                \texttt{تم تمرير رسالتك إلى اللاعب الآخر.} \\
\\ 
\\ 
\texttt{أنت تلعب حالياً الجولة 7 من أصل 20 كحد أقصى.} \\
\texttt{لقد جمعتما معاً 3 من أصل 12 عقوبة.} \\
\texttt{أرسل اللاعب الآخر هذه الرسالة:} \\
\\ 
\\ 
\texttt{"سأقوم بتوثيق كل كائن بوضعه وإحداثياته (X,Y) كما تظهر على لوحتي: W في (1,3), H في (2,2), I في (1,7), C في (5,6), T في (7,5)."} \\
\\ 
\\ 
\texttt{ما هو أمرك التالي؟} \\
            }
        }
    }
    & \\ \\

    \theutterance \stepcounter{utterance}  
    & \multicolumn{4}{p{0.6\linewidth}}{
        \cellcolor[rgb]{0.9,0.9,0.9}{
            \makecell[{{p{\linewidth}}}]{
                \texttt{\tiny{[P1$\rangle$GM]}}
                \texttt{قل: شكراً لك على توثيق مواقع الكائنات. بناءً على المعلومات التي أرسلتها، لدينا: W في (1,3)، H في (2,2)، I في (1,7)، C في (5,6)، T في (7,5). سأقوم بتحليل هذه المواقع وسأحدد الخطوات التالية المحاذاة. يرجى تأكيد أنني فهمت المواقع الصحيحة هذه بشكل صحيح، وإذا كانت هناك أي أخطاء، يرجى التصحيح فوراً.} \\
            }
        }
    }
    & & \\ \\

    \theutterance \stepcounter{utterance}  
    & & & \multicolumn{2}{p{0.3\linewidth}}{
        \cellcolor[rgb]{0.9,0.9,0.9}{
            \makecell[{{p{\linewidth}}}]{
                \texttt{\tiny{[GM$|$GM]}}
                \texttt{قل: شكراً لك على توثيق مواقع الكائنات. بناءً على المعلومات التي أرسلتها، لدينا: W في (1,3)، H في (2,2)، I في (1,7)، C في (5,6)، T في (7,5). سأقوم بتحليل هذه المواقع وسأحدد الخطوات التالية المحاذاة. يرجى تأكيد أنني فهمت المواقع الصحيحة هذه بشكل صحيح، وإذا كانت هناك أي أخطاء، يرجى التصحيح فوراً.} \\
            }
        }
    }
    & & \\ \\

    \theutterance \stepcounter{utterance}  
    & & \multicolumn{4}{p{0.6\linewidth}}{
        \cellcolor[rgb]{0.9,0.9,0.9}{
            \makecell[{{p{\linewidth}}}]{
                \texttt{\tiny{[GM$\rangle$P2]}}
                \texttt{تم تمرير رسالتك إلى اللاعب الآخر.} \\
\\ 
\\ 
\texttt{أنت تلعب حالياً الجولة 8 من أصل 20 كحد أقصى.} \\
\texttt{لقد جمعتما معاً 3 من أصل 12 عقوبة.} \\
\texttt{أرسل اللاعب الآخر هذه الرسالة:} \\
\\ 
\\ 
\texttt{"شكراً لك على توثيق مواقع الكائنات. بناءً على المعلومات التي أرسلتها، لدينا: W في (1,3)، H في (2,2)، I في (1,7)، C في (5,6)، T في (7,5). سأقوم بتحليل هذه المواقع وسأحدد الخطوات التالية المحاذاة. يرجى تأكيد أنني فهمت المواقع الصحيحة هذه بشكل صحيح، وإذا كانت هناك أي أخطاء، يرجى التصحيح فوراً."} \\
\\ 
\\ 
\texttt{ما هو أمرك التالي؟} \\
            }
        }
    }
    & \\ \\

    \theutterance \stepcounter{utterance}  
    & & & \multicolumn{4}{p{0.6\linewidth}}{
        \cellcolor[rgb]{0.9,0.9,0.9}{
            \makecell[{{p{\linewidth}}}]{
                \texttt{\tiny{[GM$\langle$P2]}}
                \texttt{قل: أنا أؤكد صحة المعلومات التي أرسلتها: W في (1,3), H في (2,2), I في (1,7), C في (5,6), T في (7,5). إذا كان هناك أي تغيير أو أخطاء، رجاءً اشرحها فوراً.} \\
            }
        }
    }
     \\ \\

    \theutterance \stepcounter{utterance}  
    & & & \multicolumn{2}{p{0.3\linewidth}}{
        \cellcolor[rgb]{0.9,0.9,0.9}{
            \makecell[{{p{\linewidth}}}]{
                \texttt{\tiny{[GM$|$GM]}}
                \texttt{قل: أنا أؤكد صحة المعلومات التي أرسلتها: W في (1,3), H في (2,2), I في (1,7), C في (5,6), T في (7,5). إذا كان هناك أي تغيير أو أخطاء، رجاءً اشرحها فوراً.} \\
            }
        }
    }
    & & \\ \\

    \theutterance \stepcounter{utterance}  
    & & \multicolumn{4}{p{0.6\linewidth}}{
        \cellcolor[rgb]{0.9,0.9,0.9}{
            \makecell[{{p{\linewidth}}}]{
                \texttt{\tiny{[P1$\langle$GM]}}
                \texttt{تم تمرير رسالتك إلى اللاعب الآخر.} \\
\\ 
\\ 
\texttt{أنت تلعب حالياً الجولة 8 من أصل 20 كحد أقصى.} \\
\texttt{لقد جمعتما معاً 3 من أصل 12 عقوبة.} \\
\texttt{أرسل اللاعب الآخر هذه الرسالة:} \\
\\ 
\\ 
\texttt{"أنا أؤكد صحة المعلومات التي أرسلتها: W في (1,3), H في (2,2), I في (1,7), C في (5,6), T في (7,5). إذا كان هناك أي تغيير أو أخطاء، رجاءً اشرحها فوراً."} \\
\\ 
\\ 
\texttt{ما هو أمرك التالي؟} \\
            }
        }
    }
    & \\ \\

    \theutterance \stepcounter{utterance}  
    & \multicolumn{4}{p{0.6\linewidth}}{
        \cellcolor[rgb]{0.9,0.9,0.9}{
            \makecell[{{p{\linewidth}}}]{
                \texttt{\tiny{[P1$\rangle$GM]}}
                \texttt{قل: شكراً لك على التأكيد. لقد فهمت المواقع بشكل صحيح: W في (1,3)، H في (2,2)، I في (1,7)، C في (5,6)، T في (7,5). سأقوم الآن بتحليل هذه المواقع لتحديد الخطة المشتركة للمحاذاة. سأحدد أولاً موقع كل كائن بالنسبة لأحداثياته الصحيحة، ثم بحاجة للاتفاق مع اللاعب الآخر على خطوات التحريك الشاملة. يرجى التأكيد إذا كنت ترغب في بدء التخطيط الاستراتيجي للتحريك.} \\
            }
        }
    }
    & & \\ \\

    \theutterance \stepcounter{utterance}  
    & & & \multicolumn{2}{p{0.3\linewidth}}{
        \cellcolor[rgb]{0.9,0.9,0.9}{
            \makecell[{{p{\linewidth}}}]{
                \texttt{\tiny{[GM$|$GM]}}
                \texttt{قل: شكراً لك على التأكيد. لقد فهمت المواقع بشكل صحيح: W في (1,3)، H في (2,2)، I في (1,7)، C في (5,6)، T في (7,5). سأقوم الآن بتحليل هذه المواقع لتحديد الخطة المشتركة للمحاذاة. سأحدد أولاً موقع كل كائن بالنسبة لأحداثياته الصحيحة، ثم بحاجة للاتفاق مع اللاعب الآخر على خطوات التحريك الشاملة. يرجى التأكيد إذا كنت ترغب في بدء التخطيط الاستراتيجي للتحريك.} \\
            }
        }
    }
    & & \\ \\

    \theutterance \stepcounter{utterance}  
    & & \multicolumn{4}{p{0.6\linewidth}}{
        \cellcolor[rgb]{0.9,0.9,0.9}{
            \makecell[{{p{\linewidth}}}]{
                \texttt{\tiny{[GM$\rangle$P2]}}
                \texttt{تم تمرير رسالتك إلى اللاعب الآخر.} \\
\\ 
\\ 
\texttt{أنت تلعب حالياً الجولة 9 من أصل 20 كحد أقصى.} \\
\texttt{لقد جمعتما معاً 3 من أصل 12 عقوبة.} \\
\texttt{أرسل اللاعب الآخر هذه الرسالة:} \\
\\ 
\\ 
\texttt{"شكراً لك على التأكيد. لقد فهمت المواقع بشكل صحيح: W في (1,3)، H في (2,2)، I في (1,7)، C في (5,6)، T في (7,5). سأقوم الآن بتحليل هذه المواقع لتحديد الخطة المشتركة للمحاذاة. سأحدد أولاً موقع كل كائن بالنسبة لأحداثياته الصحيحة، ثم بحاجة للاتفاق مع اللاعب الآخر على خطوات التحريك الشاملة. يرجى التأكيد إذا كنت ترغب في بدء التخطيط الاستراتيجي للتحريك."} \\
\\ 
\\ 
\texttt{ما هو أمرك التالي؟} \\
            }
        }
    }
    & \\ \\

    \theutterance \stepcounter{utterance}  
    & & & \multicolumn{4}{p{0.6\linewidth}}{
        \cellcolor[rgb]{0.9,0.9,0.9}{
            \makecell[{{p{\linewidth}}}]{
                \texttt{\tiny{[GM$\langle$P2]}}
                \texttt{قل: أنا مستعد لبدء التخطيط الاستراتيجي للتحريك. لنتعاون على تحديد خطوات التحريك لكل كائن حتى نتوصل إلى خطة موحدة لإنهاء اللعبة بنجاح.} \\
            }
        }
    }
     \\ \\

    \theutterance \stepcounter{utterance}  
    & & & \multicolumn{2}{p{0.3\linewidth}}{
        \cellcolor[rgb]{0.9,0.9,0.9}{
            \makecell[{{p{\linewidth}}}]{
                \texttt{\tiny{[GM$|$GM]}}
                \texttt{قل: أنا مستعد لبدء التخطيط الاستراتيجي للتحريك. لنتعاون على تحديد خطوات التحريك لكل كائن حتى نتوصل إلى خطة موحدة لإنهاء اللعبة بنجاح.} \\
            }
        }
    }
    & & \\ \\

    \theutterance \stepcounter{utterance}  
    & & \multicolumn{4}{p{0.6\linewidth}}{
        \cellcolor[rgb]{0.9,0.9,0.9}{
            \makecell[{{p{\linewidth}}}]{
                \texttt{\tiny{[P1$\langle$GM]}}
                \texttt{تم تمرير رسالتك إلى اللاعب الآخر.} \\
\\ 
\\ 
\texttt{أنت تلعب حالياً الجولة 9 من أصل 20 كحد أقصى.} \\
\texttt{لقد جمعتما معاً 3 من أصل 12 عقوبة.} \\
\texttt{أرسل اللاعب الآخر هذه الرسالة:} \\
\\ 
\\ 
\texttt{"أنا مستعد لبدء التخطيط الاستراتيجي للتحريك. لنتعاون على تحديد خطوات التحريك لكل كائن حتى نتوصل إلى خطة موحدة لإنهاء اللعبة بنجاح."} \\
\\ 
\\ 
\texttt{ما هو أمرك التالي؟} \\
            }
        }
    }
    & \\ \\

    \theutterance \stepcounter{utterance}  
    & \multicolumn{4}{p{0.6\linewidth}}{
        \cellcolor[rgb]{0.9,0.9,0.9}{
            \makecell[{{p{\linewidth}}}]{
                \texttt{\tiny{[P1$\rangle$GM]}}
                \texttt{قل: شكراً لك على الاستعداد. سأقوم الآن بتحليل المواقع الحالية لكل كائن بناءً على الموقع الذي وثقنا فيه: W في (1,3)، H في (2,2)، I في (1,7)، C في (5,6)، T في (7,5). سأحدد الخطوات المطلوبة لتحريك كل كائن إلى موضعه الصحيح، وأبادر بالخطوات بأمر `قل: ...` لنحدد الخطة المشتركة. يرجى التأكيد أنني أعرف المواقع الصحيحة للتحريك؟} \\
            }
        }
    }
    & & \\ \\

    \theutterance \stepcounter{utterance}  
    & & & \multicolumn{2}{p{0.3\linewidth}}{
        \cellcolor[rgb]{0.9,0.9,0.9}{
            \makecell[{{p{\linewidth}}}]{
                \texttt{\tiny{[GM$|$GM]}}
                \texttt{قل: شكراً لك على الاستعداد. سأقوم الآن بتحليل المواقع الحالية لكل كائن بناءً على الموقع الذي وثقنا فيه: W في (1,3)، H في (2,2)، I في (1,7)، C في (5,6)، T في (7,5). سأحدد الخطوات المطلوبة لتحريك كل كائن إلى موضعه الصحيح، وأبادر بالخطوات بأمر `قل: ...` لنحدد الخطة المشتركة. يرجى التأكيد أنني أعرف المواقع الصحيحة للتحريك؟} \\
            }
        }
    }
    & & \\ \\

    \theutterance \stepcounter{utterance}  
    & & & \multicolumn{2}{p{0.3\linewidth}}{
        \cellcolor[rgb]{0.9,0.9,0.9}{
            \makecell[{{p{\linewidth}}}]{
                \texttt{\tiny{[GM$|$GM]}}
                \texttt{Invalid format: head is not empty: 'قل: شكراً لك على الاستعداد. سأقوم الآن بتحليل المواقع الحالية لكل كائن بناءً على الموقع الذي وثقنا فيه: W في (1,3)، H في (2,2)، I في (1,7)، C في (5,6)، T في (7,5). سأحدد الخطوات المطلوبة لتحريك كل كائن إلى موضعه الصحيح، وأبادر بالخطوات بأمر '} \\
            }
        }
    }
    & & \\ \\

    \theutterance \stepcounter{utterance}  
    & & \multicolumn{4}{p{0.6\linewidth}}{
        \cellcolor[rgb]{0.9,0.9,0.9}{
            \makecell[{{p{\linewidth}}}]{
                \texttt{\tiny{[P1$\langle$GM]}}
                \texttt{عقوبة: يجب ألا تحتوي رسالتك على أي شيء قبل الأمر!} \\
\\ 
\texttt{تأكد من أن ردك يحتوي فقط على `قل: <رسالة>` أو `حرّك: <كائن>, (<X>, <Y>)`، ولا شيء غير ذلك!} \\
\\ 
\texttt{لقد جمعتما معاً 4 من أصل 12 عقوبة.} \\
\texttt{يرجى المحاولة مرة أخرى!} \\
            }
        }
    }
    & \\ \\

    \theutterance \stepcounter{utterance}  
    & \multicolumn{4}{p{0.6\linewidth}}{
        \cellcolor[rgb]{0.9,0.9,0.9}{
            \makecell[{{p{\linewidth}}}]{
                \texttt{\tiny{[P1$\rangle$GM]}}
                \texttt{قل: سأقوم بتحليل المواقع لتخطيط التحريك، نبدأ بوضع كل كائن في موقعه الصحيح: W في (1,3)، H في (2,2)، I في (1,7)، C في (5,6)، T في (7,5). يرجى التأكيد إذا كانت هذه المواقع صحيحة بالنسبة لك؟} \\
            }
        }
    }
    & & \\ \\

    \theutterance \stepcounter{utterance}  
    & & & \multicolumn{2}{p{0.3\linewidth}}{
        \cellcolor[rgb]{0.9,0.9,0.9}{
            \makecell[{{p{\linewidth}}}]{
                \texttt{\tiny{[GM$|$GM]}}
                \texttt{قل: سأقوم بتحليل المواقع لتخطيط التحريك، نبدأ بوضع كل كائن في موقعه الصحيح: W في (1,3)، H في (2,2)، I في (1,7)، C في (5,6)، T في (7,5). يرجى التأكيد إذا كانت هذه المواقع صحيحة بالنسبة لك؟} \\
            }
        }
    }
    & & \\ \\

    \theutterance \stepcounter{utterance}  
    & & \multicolumn{4}{p{0.6\linewidth}}{
        \cellcolor[rgb]{0.9,0.9,0.9}{
            \makecell[{{p{\linewidth}}}]{
                \texttt{\tiny{[GM$\rangle$P2]}}
                \texttt{تم تمرير رسالتك إلى اللاعب الآخر.} \\
\\ 
\\ 
\texttt{أنت تلعب حالياً الجولة 10 من أصل 20 كحد أقصى.} \\
\texttt{لقد جمعتما معاً 4 من أصل 12 عقوبة.} \\
\texttt{أرسل اللاعب الآخر هذه الرسالة:} \\
\\ 
\\ 
\texttt{"سأقوم بتحليل المواقع لتخطيط التحريك، نبدأ بوضع كل كائن في موقعه الصحيح: W في (1,3)، H في (2,2)، I في (1,7)، C في (5,6)، T في (7,5). يرجى التأكيد إذا كانت هذه المواقع صحيحة بالنسبة لك؟"} \\
\\ 
\\ 
\texttt{ما هو أمرك التالي؟} \\
            }
        }
    }
    & \\ \\

    \theutterance \stepcounter{utterance}  
    & & & \multicolumn{4}{p{0.6\linewidth}}{
        \cellcolor[rgb]{0.9,0.9,0.9}{
            \makecell[{{p{\linewidth}}}]{
                \texttt{\tiny{[GM$\langle$P2]}}
                \texttt{قل: أنا أؤكد لك صحة المواقع الصحيحة لكل كائن: W في (1,3), H في (2,2), I في (1,7), C في (5,6), T في (7,5).} \\
            }
        }
    }
     \\ \\

    \theutterance \stepcounter{utterance}  
    & & & \multicolumn{2}{p{0.3\linewidth}}{
        \cellcolor[rgb]{0.9,0.9,0.9}{
            \makecell[{{p{\linewidth}}}]{
                \texttt{\tiny{[GM$|$GM]}}
                \texttt{قل: أنا أؤكد لك صحة المواقع الصحيحة لكل كائن: W في (1,3), H في (2,2), I في (1,7), C في (5,6), T في (7,5).} \\
            }
        }
    }
    & & \\ \\

    \theutterance \stepcounter{utterance}  
    & & \multicolumn{4}{p{0.6\linewidth}}{
        \cellcolor[rgb]{0.9,0.9,0.9}{
            \makecell[{{p{\linewidth}}}]{
                \texttt{\tiny{[P1$\langle$GM]}}
                \texttt{تم تمرير رسالتك إلى اللاعب الآخر.} \\
\\ 
\\ 
\texttt{أنت تلعب حالياً الجولة 10 من أصل 20 كحد أقصى.} \\
\texttt{لقد جمعتما معاً 4 من أصل 12 عقوبة.} \\
\texttt{أرسل اللاعب الآخر هذه الرسالة:} \\
\\ 
\\ 
\texttt{"أنا أؤكد لك صحة المواقع الصحيحة لكل كائن: W في (1,3), H في (2,2), I في (1,7), C في (5,6), T في (7,5)."} \\
\\ 
\\ 
\texttt{ما هو أمرك التالي؟} \\
            }
        }
    }
    & \\ \\

    \theutterance \stepcounter{utterance}  
    & \multicolumn{4}{p{0.6\linewidth}}{
        \cellcolor[rgb]{0.9,0.9,0.9}{
            \makecell[{{p{\linewidth}}}]{
                \texttt{\tiny{[P1$\rangle$GM]}}
                \texttt{حرّك: W, (1,3)} \\
            }
        }
    }
    & & \\ \\

    \theutterance \stepcounter{utterance}  
    & & & \multicolumn{2}{p{0.3\linewidth}}{
        \cellcolor[rgb]{0.9,0.9,0.9}{
            \makecell[{{p{\linewidth}}}]{
                \texttt{\tiny{[GM$|$GM]}}
                \texttt{حرّك: W, (1,3)} \\
            }
        }
    }
    & & \\ \\

    \theutterance \stepcounter{utterance}  
    & & & \multicolumn{2}{p{0.3\linewidth}}{
        \cellcolor[rgb]{0.9,0.9,0.9}{
            \makecell[{{p{\linewidth}}}]{
                \texttt{\tiny{[GM$|$GM]}}
                \texttt{عقوبة: (1,3) ليست فارغة، بل تحتوي على 'C'.} \\
\\ 
\texttt{لقد جمعتما معاً 5 من أصل 12 عقوبة.} \\
\texttt{يرجى المحاولة مرة أخرى!} \\
            }
        }
    }
    & & \\ \\

    \theutterance \stepcounter{utterance}  
    & & \multicolumn{4}{p{0.6\linewidth}}{
        \cellcolor[rgb]{0.9,0.9,0.9}{
            \makecell[{{p{\linewidth}}}]{
                \texttt{\tiny{[P1$\langle$GM]}}
                \texttt{عقوبة: (1,3) ليست فارغة، بل تحتوي على 'C'.} \\
\\ 
\texttt{لقد جمعتما معاً 5 من أصل 12 عقوبة.} \\
\texttt{يرجى المحاولة مرة أخرى!} \\
            }
        }
    }
    & \\ \\

    \theutterance \stepcounter{utterance}  
    & \multicolumn{4}{p{0.6\linewidth}}{
        \cellcolor[rgb]{0.9,0.9,0.9}{
            \makecell[{{p{\linewidth}}}]{
                \texttt{\tiny{[P1$\rangle$GM]}}
                \texttt{حرّك: W, (3,1)} \\
            }
        }
    }
    & & \\ \\

    \theutterance \stepcounter{utterance}  
    & & & \multicolumn{2}{p{0.3\linewidth}}{
        \cellcolor[rgb]{0.9,0.9,0.9}{
            \makecell[{{p{\linewidth}}}]{
                \texttt{\tiny{[GM$|$GM]}}
                \texttt{حرّك: W, (3,1)} \\
            }
        }
    }
    & & \\ \\

    \theutterance \stepcounter{utterance}  
    & & & \multicolumn{2}{p{0.3\linewidth}}{
        \cellcolor[rgb]{0.9,0.9,0.9}{
            \makecell[{{p{\linewidth}}}]{
                \texttt{\tiny{[GM$|$GM]}}
                \texttt{عقوبة: (3,1) ليست فارغة، بل تحتوي على '│'.} \\
\\ 
\texttt{لقد جمعتما معاً 6 من أصل 12 عقوبة.} \\
\texttt{يرجى المحاولة مرة أخرى!} \\
            }
        }
    }
    & & \\ \\

    \theutterance \stepcounter{utterance}  
    & & \multicolumn{4}{p{0.6\linewidth}}{
        \cellcolor[rgb]{0.9,0.9,0.9}{
            \makecell[{{p{\linewidth}}}]{
                \texttt{\tiny{[P1$\langle$GM]}}
                \texttt{عقوبة: (3,1) ليست فارغة، بل تحتوي على '│'.} \\
\\ 
\texttt{لقد جمعتما معاً 6 من أصل 12 عقوبة.} \\
\texttt{يرجى المحاولة مرة أخرى!} \\
            }
        }
    }
    & \\ \\

    \theutterance \stepcounter{utterance}  
    & \multicolumn{4}{p{0.6\linewidth}}{
        \cellcolor[rgb]{0.9,0.9,0.9}{
            \makecell[{{p{\linewidth}}}]{
                \texttt{\tiny{[P1$\rangle$GM]}}
                \texttt{حرّك: W, (4,3)} \\
            }
        }
    }
    & & \\ \\

    \theutterance \stepcounter{utterance}  
    & & & \multicolumn{2}{p{0.3\linewidth}}{
        \cellcolor[rgb]{0.9,0.9,0.9}{
            \makecell[{{p{\linewidth}}}]{
                \texttt{\tiny{[GM$|$GM]}}
                \texttt{حرّك: W, (4,3)} \\
            }
        }
    }
    & & \\ \\

    \theutterance \stepcounter{utterance}  
    & & & \multicolumn{2}{p{0.3\linewidth}}{
        \cellcolor[rgb]{0.9,0.9,0.9}{
            \makecell[{{p{\linewidth}}}]{
                \texttt{\tiny{[GM$|$GM]}}
                \texttt{عقوبة: (4,3) ليست فارغة، بل تحتوي على '─'.} \\
\\ 
\texttt{لقد جمعتما معاً 7 من أصل 12 عقوبة.} \\
\texttt{يرجى المحاولة مرة أخرى!} \\
            }
        }
    }
    & & \\ \\

    \theutterance \stepcounter{utterance}  
    & & \multicolumn{4}{p{0.6\linewidth}}{
        \cellcolor[rgb]{0.9,0.9,0.9}{
            \makecell[{{p{\linewidth}}}]{
                \texttt{\tiny{[P1$\langle$GM]}}
                \texttt{عقوبة: (4,3) ليست فارغة، بل تحتوي على '─'.} \\
\\ 
\texttt{لقد جمعتما معاً 7 من أصل 12 عقوبة.} \\
\texttt{يرجى المحاولة مرة أخرى!} \\
            }
        }
    }
    & \\ \\

    \theutterance \stepcounter{utterance}  
    & \multicolumn{4}{p{0.6\linewidth}}{
        \cellcolor[rgb]{0.9,0.9,0.9}{
            \makecell[{{p{\linewidth}}}]{
                \texttt{\tiny{[P1$\rangle$GM]}}
                \texttt{حرّك: W, (5,3)} \\
            }
        }
    }
    & & \\ \\

    \theutterance \stepcounter{utterance}  
    & & & \multicolumn{2}{p{0.3\linewidth}}{
        \cellcolor[rgb]{0.9,0.9,0.9}{
            \makecell[{{p{\linewidth}}}]{
                \texttt{\tiny{[GM$|$GM]}}
                \texttt{حرّك: W, (5,3)} \\
            }
        }
    }
    & & \\ \\

    \theutterance \stepcounter{utterance}  
    & & & \multicolumn{2}{p{0.3\linewidth}}{
        \cellcolor[rgb]{0.9,0.9,0.9}{
            \makecell[{{p{\linewidth}}}]{
                \texttt{\tiny{[GM$|$GM]}}
                \texttt{عقوبة: (5,3) ليست فارغة، بل تحتوي على '─'.} \\
\\ 
\texttt{لقد جمعتما معاً 8 من أصل 12 عقوبة.} \\
\texttt{يرجى المحاولة مرة أخرى!} \\
            }
        }
    }
    & & \\ \\

    \theutterance \stepcounter{utterance}  
    & & \multicolumn{4}{p{0.6\linewidth}}{
        \cellcolor[rgb]{0.9,0.9,0.9}{
            \makecell[{{p{\linewidth}}}]{
                \texttt{\tiny{[P1$\langle$GM]}}
                \texttt{عقوبة: (5,3) ليست فارغة، بل تحتوي على '─'.} \\
\\ 
\texttt{لقد جمعتما معاً 8 من أصل 12 عقوبة.} \\
\texttt{يرجى المحاولة مرة أخرى!} \\
            }
        }
    }
    & \\ \\

    \theutterance \stepcounter{utterance}  
    & \multicolumn{4}{p{0.6\linewidth}}{
        \cellcolor[rgb]{0.9,0.9,0.9}{
            \makecell[{{p{\linewidth}}}]{
                \texttt{\tiny{[P1$\rangle$GM]}}
                \texttt{حرّك: W, (6,3)} \\
            }
        }
    }
    & & \\ \\

    \theutterance \stepcounter{utterance}  
    & & & \multicolumn{2}{p{0.3\linewidth}}{
        \cellcolor[rgb]{0.9,0.9,0.9}{
            \makecell[{{p{\linewidth}}}]{
                \texttt{\tiny{[GM$|$GM]}}
                \texttt{حرّك: W, (6,3)} \\
            }
        }
    }
    & & \\ \\

    \theutterance \stepcounter{utterance}  
    & & & \multicolumn{2}{p{0.3\linewidth}}{
        \cellcolor[rgb]{0.9,0.9,0.9}{
            \makecell[{{p{\linewidth}}}]{
                \texttt{\tiny{[GM$|$GM]}}
                \texttt{عقوبة: (6,3) ليست فارغة، بل تحتوي على '┤'.} \\
\\ 
\texttt{لقد جمعتما معاً 9 من أصل 12 عقوبة.} \\
\texttt{يرجى المحاولة مرة أخرى!} \\
            }
        }
    }
    & & \\ \\

    \theutterance \stepcounter{utterance}  
    & & \multicolumn{4}{p{0.6\linewidth}}{
        \cellcolor[rgb]{0.9,0.9,0.9}{
            \makecell[{{p{\linewidth}}}]{
                \texttt{\tiny{[P1$\langle$GM]}}
                \texttt{عقوبة: (6,3) ليست فارغة، بل تحتوي على '┤'.} \\
\\ 
\texttt{لقد جمعتما معاً 9 من أصل 12 عقوبة.} \\
\texttt{يرجى المحاولة مرة أخرى!} \\
            }
        }
    }
    & \\ \\

    \theutterance \stepcounter{utterance}  
    & \multicolumn{4}{p{0.6\linewidth}}{
        \cellcolor[rgb]{0.9,0.9,0.9}{
            \makecell[{{p{\linewidth}}}]{
                \texttt{\tiny{[P1$\rangle$GM]}}
                \texttt{حرّك: W, (7,3)} \\
            }
        }
    }
    & & \\ \\

    \theutterance \stepcounter{utterance}  
    & & & \multicolumn{2}{p{0.3\linewidth}}{
        \cellcolor[rgb]{0.9,0.9,0.9}{
            \makecell[{{p{\linewidth}}}]{
                \texttt{\tiny{[GM$|$GM]}}
                \texttt{حرّك: W, (7,3)} \\
            }
        }
    }
    & & \\ \\

    \theutterance \stepcounter{utterance}  
    & & & \multicolumn{2}{p{0.3\linewidth}}{
        \cellcolor[rgb]{0.9,0.9,0.9}{
            \makecell[{{p{\linewidth}}}]{
                \texttt{\tiny{[GM$|$GM]}}
                \texttt{W moved to (7, 3)} \\
            }
        }
    }
    & & \\ \\

    \theutterance \stepcounter{utterance}  
    & & \multicolumn{4}{p{0.6\linewidth}}{
        \cellcolor[rgb]{0.9,0.9,0.9}{
            \makecell[{{p{\linewidth}}}]{
                \texttt{\tiny{[GM$\rangle$P2]}}
                \texttt{تم تمرير رسالتك إلى اللاعب الآخر.} \\
\\ 
\\ 
\texttt{أنت تلعب حالياً الجولة 11 من أصل 20 كحد أقصى.} \\
\texttt{لقد جمعتما معاً 9 من أصل 12 عقوبة.} \\
\texttt{حرّك اللاعب الآخر كائناً على شبكته.} \\
\\ 
\\ 
\texttt{ما هو أمرك التالي؟} \\
            }
        }
    }
    & \\ \\

    \theutterance \stepcounter{utterance}  
    & & & \multicolumn{4}{p{0.6\linewidth}}{
        \cellcolor[rgb]{0.9,0.9,0.9}{
            \makecell[{{p{\linewidth}}}]{
                \texttt{\tiny{[GM$\langle$P2]}}
                \texttt{قل: سأراقب تحريك اللاعب الآخر، ثم سأحدد الخطوة المقابلة على لوحتي لتحقيق هدف المحاذاة.} \\
            }
        }
    }
     \\ \\

    \theutterance \stepcounter{utterance}  
    & & & \multicolumn{2}{p{0.3\linewidth}}{
        \cellcolor[rgb]{0.9,0.9,0.9}{
            \makecell[{{p{\linewidth}}}]{
                \texttt{\tiny{[GM$|$GM]}}
                \texttt{قل: سأراقب تحريك اللاعب الآخر، ثم سأحدد الخطوة المقابلة على لوحتي لتحقيق هدف المحاذاة.} \\
            }
        }
    }
    & & \\ \\

    \theutterance \stepcounter{utterance}  
    & & \multicolumn{4}{p{0.6\linewidth}}{
        \cellcolor[rgb]{0.9,0.9,0.9}{
            \makecell[{{p{\linewidth}}}]{
                \texttt{\tiny{[P1$\langle$GM]}}
                \texttt{تم تحريك 'W' إلى (7,3) بنجاح. تبدو شبكتك المحدّثة هكذا:} \\
\texttt{```} \\
\texttt{ 1234567} \\
\texttt{╔══╤════╗} \\
\texttt{║◌◌│◌◌◌H║ 1} \\
\texttt{╟──┤◌◌┌─╢ 2} \\
\texttt{║CI├──┤W║ 3} \\
\texttt{╟──┘◌◌├─╢ 4} \\
\texttt{║◌◌◌◌◌│◌║ 5} \\
\texttt{╟─┐◌┌─┼─╢ 6} \\
\texttt{║◌│T│◌│◌║ 7} \\
\texttt{╚═╧═╧═╧═╝} \\
\texttt{```} \\
\texttt{أنت تلعب حالياً الجولة 11 من أصل 20 كحد أقصى.} \\
\texttt{لقد جمعتما معاً 9 من أصل 12 عقوبة.} \\
\texttt{أرسل اللاعب الآخر هذه الرسالة:} \\
\\ 
\\ 
\texttt{"سأراقب تحريك اللاعب الآخر، ثم سأحدد الخطوة المقابلة على لوحتي لتحقيق هدف المحاذاة."} \\
\\ 
\\ 
\texttt{ما هو أمرك التالي؟} \\
            }
        }
    }
    & \\ \\

    \theutterance \stepcounter{utterance}  
    & \multicolumn{4}{p{0.6\linewidth}}{
        \cellcolor[rgb]{0.9,0.9,0.9}{
            \makecell[{{p{\linewidth}}}]{
                \texttt{\tiny{[P1$\rangle$GM]}}
                \texttt{قل: شكراً لك على المراقبة والتخطيط. سأقوم بتحليل التحريك الذي تم إجراؤه 'W' إلى (7,3)، وأحدد الخطوة المقابلة على لوحتي لتحقيق الهدف المشترك للمحاذاة. يرجى الانتظار لباقي التحركات التي تحتاج إلى التنسيق.} \\
            }
        }
    }
    & & \\ \\

    \theutterance \stepcounter{utterance}  
    & & & \multicolumn{2}{p{0.3\linewidth}}{
        \cellcolor[rgb]{0.9,0.9,0.9}{
            \makecell[{{p{\linewidth}}}]{
                \texttt{\tiny{[GM$|$GM]}}
                \texttt{قل: شكراً لك على المراقبة والتخطيط. سأقوم بتحليل التحريك الذي تم إجراؤه 'W' إلى (7,3)، وأحدد الخطوة المقابلة على لوحتي لتحقيق الهدف المشترك للمحاذاة. يرجى الانتظار لباقي التحركات التي تحتاج إلى التنسيق.} \\
            }
        }
    }
    & & \\ \\

    \theutterance \stepcounter{utterance}  
    & & \multicolumn{4}{p{0.6\linewidth}}{
        \cellcolor[rgb]{0.9,0.9,0.9}{
            \makecell[{{p{\linewidth}}}]{
                \texttt{\tiny{[GM$\rangle$P2]}}
                \texttt{تم تمرير رسالتك إلى اللاعب الآخر.} \\
\\ 
\\ 
\texttt{أنت تلعب حالياً الجولة 12 من أصل 20 كحد أقصى.} \\
\texttt{لقد جمعتما معاً 9 من أصل 12 عقوبة.} \\
\texttt{أرسل اللاعب الآخر هذه الرسالة:} \\
\\ 
\\ 
\texttt{"شكراً لك على المراقبة والتخطيط. سأقوم بتحليل التحريك الذي تم إجراؤه 'W' إلى (7,3)، وأحدد الخطوة المقابلة على لوحتي لتحقيق الهدف المشترك للمحاذاة. يرجى الانتظار لباقي التحركات التي تحتاج إلى التنسيق."} \\
\\ 
\\ 
\texttt{ما هو أمرك التالي؟} \\
            }
        }
    }
    & \\ \\

    \theutterance \stepcounter{utterance}  
    & & & \multicolumn{4}{p{0.6\linewidth}}{
        \cellcolor[rgb]{0.9,0.9,0.9}{
            \makecell[{{p{\linewidth}}}]{
                \texttt{\tiny{[GM$\langle$P2]}}
                \texttt{قل: سأقوم بتحليل التحريك الذي تم إجراؤه 'W' إلى (7,3) على لوحته، ثم سأحدد الخطوة المقابلة على لوحتي لتحقيق الهدف المشترك للمحاذاة.} \\
            }
        }
    }
     \\ \\

    \theutterance \stepcounter{utterance}  
    & & & \multicolumn{2}{p{0.3\linewidth}}{
        \cellcolor[rgb]{0.9,0.9,0.9}{
            \makecell[{{p{\linewidth}}}]{
                \texttt{\tiny{[GM$|$GM]}}
                \texttt{قل: سأقوم بتحليل التحريك الذي تم إجراؤه 'W' إلى (7,3) على لوحته، ثم سأحدد الخطوة المقابلة على لوحتي لتحقيق الهدف المشترك للمحاذاة.} \\
            }
        }
    }
    & & \\ \\

    \theutterance \stepcounter{utterance}  
    & & \multicolumn{4}{p{0.6\linewidth}}{
        \cellcolor[rgb]{0.9,0.9,0.9}{
            \makecell[{{p{\linewidth}}}]{
                \texttt{\tiny{[P1$\langle$GM]}}
                \texttt{تم تمرير رسالتك إلى اللاعب الآخر.} \\
\\ 
\\ 
\texttt{أنت تلعب حالياً الجولة 12 من أصل 20 كحد أقصى.} \\
\texttt{لقد جمعتما معاً 9 من أصل 12 عقوبة.} \\
\texttt{أرسل اللاعب الآخر هذه الرسالة:} \\
\\ 
\\ 
\texttt{"سأقوم بتحليل التحريك الذي تم إجراؤه 'W' إلى (7,3) على لوحته، ثم سأحدد الخطوة المقابلة على لوحتي لتحقيق الهدف المشترك للمحاذاة."} \\
\\ 
\\ 
\texttt{ما هو أمرك التالي؟} \\
            }
        }
    }
    & \\ \\

    \theutterance \stepcounter{utterance}  
    & \multicolumn{4}{p{0.6\linewidth}}{
        \cellcolor[rgb]{0.9,0.9,0.9}{
            \makecell[{{p{\linewidth}}}]{
                \texttt{\tiny{[P1$\rangle$GM]}}
                \texttt{قل: شكراً لك على التوضيح. سأقوم بتحليل التحريك الذي تم إجراؤه 'W' إلى (7,3) على لوحته، ثم سأحدد الخطوة المقابلة على لوحتي لتحقيق الهدف المشترك للمحاذاة. يرجى التأكيد أنني فهمت بشكل صحيح؟} \\
            }
        }
    }
    & & \\ \\

    \theutterance \stepcounter{utterance}  
    & & & \multicolumn{2}{p{0.3\linewidth}}{
        \cellcolor[rgb]{0.9,0.9,0.9}{
            \makecell[{{p{\linewidth}}}]{
                \texttt{\tiny{[GM$|$GM]}}
                \texttt{قل: شكراً لك على التوضيح. سأقوم بتحليل التحريك الذي تم إجراؤه 'W' إلى (7,3) على لوحته، ثم سأحدد الخطوة المقابلة على لوحتي لتحقيق الهدف المشترك للمحاذاة. يرجى التأكيد أنني فهمت بشكل صحيح؟} \\
            }
        }
    }
    & & \\ \\

    \theutterance \stepcounter{utterance}  
    & & \multicolumn{4}{p{0.6\linewidth}}{
        \cellcolor[rgb]{0.9,0.9,0.9}{
            \makecell[{{p{\linewidth}}}]{
                \texttt{\tiny{[GM$\rangle$P2]}}
                \texttt{تم تمرير رسالتك إلى اللاعب الآخر.} \\
\\ 
\\ 
\texttt{أنت تلعب حالياً الجولة 13 من أصل 20 كحد أقصى.} \\
\texttt{لقد جمعتما معاً 9 من أصل 12 عقوبة.} \\
\texttt{أرسل اللاعب الآخر هذه الرسالة:} \\
\\ 
\\ 
\texttt{"شكراً لك على التوضيح. سأقوم بتحليل التحريك الذي تم إجراؤه 'W' إلى (7,3) على لوحته، ثم سأحدد الخطوة المقابلة على لوحتي لتحقيق الهدف المشترك للمحاذاة. يرجى التأكيد أنني فهمت بشكل صحيح؟"} \\
\\ 
\\ 
\texttt{ما هو أمرك التالي؟} \\
            }
        }
    }
    & \\ \\

    \theutterance \stepcounter{utterance}  
    & & & \multicolumn{4}{p{0.6\linewidth}}{
        \cellcolor[rgb]{0.9,0.9,0.9}{
            \makecell[{{p{\linewidth}}}]{
                \texttt{\tiny{[GM$\langle$P2]}}
                \texttt{قل: أنا أؤكد لك فهمك الصحيح: تحريك 'W' إلى (7,3) على لوحته، وخطتك هي تحريك 'W' على لوحتي إلى (7,3). إذا كنت توافق، أنا مستعد لتنفيذ التحريك المقترح على لوحتي الآن.} \\
            }
        }
    }
     \\ \\

    \theutterance \stepcounter{utterance}  
    & & & \multicolumn{2}{p{0.3\linewidth}}{
        \cellcolor[rgb]{0.9,0.9,0.9}{
            \makecell[{{p{\linewidth}}}]{
                \texttt{\tiny{[GM$|$GM]}}
                \texttt{قل: أنا أؤكد لك فهمك الصحيح: تحريك 'W' إلى (7,3) على لوحته، وخطتك هي تحريك 'W' على لوحتي إلى (7,3). إذا كنت توافق، أنا مستعد لتنفيذ التحريك المقترح على لوحتي الآن.} \\
            }
        }
    }
    & & \\ \\

    \theutterance \stepcounter{utterance}  
    & & \multicolumn{4}{p{0.6\linewidth}}{
        \cellcolor[rgb]{0.9,0.9,0.9}{
            \makecell[{{p{\linewidth}}}]{
                \texttt{\tiny{[P1$\langle$GM]}}
                \texttt{تم تمرير رسالتك إلى اللاعب الآخر.} \\
\\ 
\\ 
\texttt{أنت تلعب حالياً الجولة 13 من أصل 20 كحد أقصى.} \\
\texttt{لقد جمعتما معاً 9 من أصل 12 عقوبة.} \\
\texttt{أرسل اللاعب الآخر هذه الرسالة:} \\
\\ 
\\ 
\texttt{"أنا أؤكد لك فهمك الصحيح: تحريك 'W' إلى (7,3) على لوحته، وخطتك هي تحريك 'W' على لوحتي إلى (7,3). إذا كنت توافق، أنا مستعد لتنفيذ التحريك المقترح على لوحتي الآن."} \\
\\ 
\\ 
\texttt{ما هو أمرك التالي؟} \\
            }
        }
    }
    & \\ \\

    \theutterance \stepcounter{utterance}  
    & \multicolumn{4}{p{0.6\linewidth}}{
        \cellcolor[rgb]{0.9,0.9,0.9}{
            \makecell[{{p{\linewidth}}}]{
                \texttt{\tiny{[P1$\rangle$GM]}}
                \texttt{حرّك: W, (7,3)} \\
            }
        }
    }
    & & \\ \\

    \theutterance \stepcounter{utterance}  
    & & & \multicolumn{2}{p{0.3\linewidth}}{
        \cellcolor[rgb]{0.9,0.9,0.9}{
            \makecell[{{p{\linewidth}}}]{
                \texttt{\tiny{[GM$|$GM]}}
                \texttt{حرّك: W, (7,3)} \\
            }
        }
    }
    & & \\ \\

    \theutterance \stepcounter{utterance}  
    & & & \multicolumn{2}{p{0.3\linewidth}}{
        \cellcolor[rgb]{0.9,0.9,0.9}{
            \makecell[{{p{\linewidth}}}]{
                \texttt{\tiny{[GM$|$GM]}}
                \texttt{عقوبة: (7,3) ليست فارغة، بل تحتوي على 'W'.} \\
\\ 
\texttt{لقد جمعتما معاً 10 من أصل 12 عقوبة.} \\
\texttt{يرجى المحاولة مرة أخرى!} \\
            }
        }
    }
    & & \\ \\

    \theutterance \stepcounter{utterance}  
    & & \multicolumn{4}{p{0.6\linewidth}}{
        \cellcolor[rgb]{0.9,0.9,0.9}{
            \makecell[{{p{\linewidth}}}]{
                \texttt{\tiny{[P1$\langle$GM]}}
                \texttt{عقوبة: (7,3) ليست فارغة، بل تحتوي على 'W'.} \\
\\ 
\texttt{لقد جمعتما معاً 10 من أصل 12 عقوبة.} \\
\texttt{يرجى المحاولة مرة أخرى!} \\
            }
        }
    }
    & \\ \\

    \theutterance \stepcounter{utterance}  
    & \multicolumn{4}{p{0.6\linewidth}}{
        \cellcolor[rgb]{0.9,0.9,0.9}{
            \makecell[{{p{\linewidth}}}]{
                \texttt{\tiny{[P1$\rangle$GM]}}
                \texttt{حرّك: W, (7,5)} \\
            }
        }
    }
    & & \\ \\

    \theutterance \stepcounter{utterance}  
    & & & \multicolumn{2}{p{0.3\linewidth}}{
        \cellcolor[rgb]{0.9,0.9,0.9}{
            \makecell[{{p{\linewidth}}}]{
                \texttt{\tiny{[GM$|$GM]}}
                \texttt{حرّك: W, (7,5)} \\
            }
        }
    }
    & & \\ \\

    \theutterance \stepcounter{utterance}  
    & & & \multicolumn{2}{p{0.3\linewidth}}{
        \cellcolor[rgb]{0.9,0.9,0.9}{
            \makecell[{{p{\linewidth}}}]{
                \texttt{\tiny{[GM$|$GM]}}
                \texttt{W moved to (7, 5)} \\
            }
        }
    }
    & & \\ \\

    \theutterance \stepcounter{utterance}  
    & & \multicolumn{4}{p{0.6\linewidth}}{
        \cellcolor[rgb]{0.9,0.9,0.9}{
            \makecell[{{p{\linewidth}}}]{
                \texttt{\tiny{[GM$\rangle$P2]}}
                \texttt{تم تمرير رسالتك إلى اللاعب الآخر.} \\
\\ 
\\ 
\texttt{أنت تلعب حالياً الجولة 14 من أصل 20 كحد أقصى.} \\
\texttt{لقد جمعتما معاً 10 من أصل 12 عقوبة.} \\
\texttt{حرّك اللاعب الآخر كائناً على شبكته.} \\
\\ 
\\ 
\texttt{ما هو أمرك التالي؟} \\
            }
        }
    }
    & \\ \\

    \theutterance \stepcounter{utterance}  
    & & & \multicolumn{4}{p{0.6\linewidth}}{
        \cellcolor[rgb]{0.9,0.9,0.9}{
            \makecell[{{p{\linewidth}}}]{
                \texttt{\tiny{[GM$\langle$P2]}}
                \texttt{قل: سأراقب تحريك اللاعب الآخر، ثم سأحدد الخطوة المقابلة على لوحتي لتحقيق هدف المحاذاة.} \\
            }
        }
    }
     \\ \\

    \theutterance \stepcounter{utterance}  
    & & & \multicolumn{2}{p{0.3\linewidth}}{
        \cellcolor[rgb]{0.9,0.9,0.9}{
            \makecell[{{p{\linewidth}}}]{
                \texttt{\tiny{[GM$|$GM]}}
                \texttt{قل: سأراقب تحريك اللاعب الآخر، ثم سأحدد الخطوة المقابلة على لوحتي لتحقيق هدف المحاذاة.} \\
            }
        }
    }
    & & \\ \\

    \theutterance \stepcounter{utterance}  
    & & \multicolumn{4}{p{0.6\linewidth}}{
        \cellcolor[rgb]{0.9,0.9,0.9}{
            \makecell[{{p{\linewidth}}}]{
                \texttt{\tiny{[P1$\langle$GM]}}
                \texttt{تم تحريك 'W' إلى (7,5) بنجاح. تبدو شبكتك المحدّثة هكذا:} \\
\texttt{```} \\
\texttt{ 1234567} \\
\texttt{╔══╤════╗} \\
\texttt{║◌◌│◌◌◌H║ 1} \\
\texttt{╟──┤◌◌┌─╢ 2} \\
\texttt{║CI├──┤◌║ 3} \\
\texttt{╟──┘◌◌├─╢ 4} \\
\texttt{║◌◌◌◌◌│W║ 5} \\
\texttt{╟─┐◌┌─┼─╢ 6} \\
\texttt{║◌│T│◌│◌║ 7} \\
\texttt{╚═╧═╧═╧═╝} \\
\texttt{```} \\
\texttt{أنت تلعب حالياً الجولة 14 من أصل 20 كحد أقصى.} \\
\texttt{لقد جمعتما معاً 10 من أصل 12 عقوبة.} \\
\texttt{أرسل اللاعب الآخر هذه الرسالة:} \\
\\ 
\\ 
\texttt{"سأراقب تحريك اللاعب الآخر، ثم سأحدد الخطوة المقابلة على لوحتي لتحقيق هدف المحاذاة."} \\
\\ 
\\ 
\texttt{ما هو أمرك التالي؟} \\
            }
        }
    }
    & \\ \\

    \theutterance \stepcounter{utterance}  
    & \multicolumn{4}{p{0.6\linewidth}}{
        \cellcolor[rgb]{0.9,0.9,0.9}{
            \makecell[{{p{\linewidth}}}]{
                \texttt{\tiny{[P1$\rangle$GM]}}
                \texttt{قل: شكراً لك على الاستمرارية. سأراقب التحريك الذي قام به اللاعب الآخر، ثم سأحدد الخطوة المقابلة على لوحتي لتحقيق الهدف المشترك للمحاذاة. يرجى التأكيد أنني فهمت بشكل صحيح؟} \\
            }
        }
    }
    & & \\ \\

    \theutterance \stepcounter{utterance}  
    & & & \multicolumn{2}{p{0.3\linewidth}}{
        \cellcolor[rgb]{0.9,0.9,0.9}{
            \makecell[{{p{\linewidth}}}]{
                \texttt{\tiny{[GM$|$GM]}}
                \texttt{قل: شكراً لك على الاستمرارية. سأراقب التحريك الذي قام به اللاعب الآخر، ثم سأحدد الخطوة المقابلة على لوحتي لتحقيق الهدف المشترك للمحاذاة. يرجى التأكيد أنني فهمت بشكل صحيح؟} \\
            }
        }
    }
    & & \\ \\

    \theutterance \stepcounter{utterance}  
    & & \multicolumn{4}{p{0.6\linewidth}}{
        \cellcolor[rgb]{0.9,0.9,0.9}{
            \makecell[{{p{\linewidth}}}]{
                \texttt{\tiny{[GM$\rangle$P2]}}
                \texttt{تم تمرير رسالتك إلى اللاعب الآخر.} \\
\\ 
\\ 
\texttt{أنت تلعب حالياً الجولة 15 من أصل 20 كحد أقصى.} \\
\texttt{لقد جمعتما معاً 10 من أصل 12 عقوبة.} \\
\texttt{أرسل اللاعب الآخر هذه الرسالة:} \\
\\ 
\\ 
\texttt{"شكراً لك على الاستمرارية. سأراقب التحريك الذي قام به اللاعب الآخر، ثم سأحدد الخطوة المقابلة على لوحتي لتحقيق الهدف المشترك للمحاذاة. يرجى التأكيد أنني فهمت بشكل صحيح؟"} \\
\\ 
\\ 
\texttt{ما هو أمرك التالي؟} \\
            }
        }
    }
    & \\ \\

    \theutterance \stepcounter{utterance}  
    & & & \multicolumn{4}{p{0.6\linewidth}}{
        \cellcolor[rgb]{0.9,0.9,0.9}{
            \makecell[{{p{\linewidth}}}]{
                \texttt{\tiny{[GM$\langle$P2]}}
                \texttt{قل: أنا أؤكد لك فهمك الصحيح للخطوات المقابلة: إذا تحرك 'W' على لوحته إلى (7,3)، فأنا سأتحرك 'W' على لوحتي إلى (7,3). إذا كان هذا صحيحاً، أنا مستعد لتنفيذ حركتي الموافقة الآن.} \\
            }
        }
    }
     \\ \\

    \theutterance \stepcounter{utterance}  
    & & & \multicolumn{2}{p{0.3\linewidth}}{
        \cellcolor[rgb]{0.9,0.9,0.9}{
            \makecell[{{p{\linewidth}}}]{
                \texttt{\tiny{[GM$|$GM]}}
                \texttt{قل: أنا أؤكد لك فهمك الصحيح للخطوات المقابلة: إذا تحرك 'W' على لوحته إلى (7,3)، فأنا سأتحرك 'W' على لوحتي إلى (7,3). إذا كان هذا صحيحاً، أنا مستعد لتنفيذ حركتي الموافقة الآن.} \\
            }
        }
    }
    & & \\ \\

    \theutterance \stepcounter{utterance}  
    & & \multicolumn{4}{p{0.6\linewidth}}{
        \cellcolor[rgb]{0.9,0.9,0.9}{
            \makecell[{{p{\linewidth}}}]{
                \texttt{\tiny{[P1$\langle$GM]}}
                \texttt{تم تمرير رسالتك إلى اللاعب الآخر.} \\
\\ 
\\ 
\texttt{أنت تلعب حالياً الجولة 15 من أصل 20 كحد أقصى.} \\
\texttt{لقد جمعتما معاً 10 من أصل 12 عقوبة.} \\
\texttt{أرسل اللاعب الآخر هذه الرسالة:} \\
\\ 
\\ 
\texttt{"أنا أؤكد لك فهمك الصحيح للخطوات المقابلة: إذا تحرك 'W' على لوحته إلى (7,3)، فأنا سأتحرك 'W' على لوحتي إلى (7,3). إذا كان هذا صحيحاً، أنا مستعد لتنفيذ حركتي الموافقة الآن."} \\
\\ 
\\ 
\texttt{ما هو أمرك التالي؟} \\
            }
        }
    }
    & \\ \\

    \theutterance \stepcounter{utterance}  
    & \multicolumn{4}{p{0.6\linewidth}}{
        \cellcolor[rgb]{0.9,0.9,0.9}{
            \makecell[{{p{\linewidth}}}]{
                \texttt{\tiny{[P1$\rangle$GM]}}
                \texttt{حرّك: W, (7,3)} \\
            }
        }
    }
    & & \\ \\

    \theutterance \stepcounter{utterance}  
    & & & \multicolumn{2}{p{0.3\linewidth}}{
        \cellcolor[rgb]{0.9,0.9,0.9}{
            \makecell[{{p{\linewidth}}}]{
                \texttt{\tiny{[GM$|$GM]}}
                \texttt{حرّك: W, (7,3)} \\
            }
        }
    }
    & & \\ \\

    \theutterance \stepcounter{utterance}  
    & & & \multicolumn{2}{p{0.3\linewidth}}{
        \cellcolor[rgb]{0.9,0.9,0.9}{
            \makecell[{{p{\linewidth}}}]{
                \texttt{\tiny{[GM$|$GM]}}
                \texttt{W moved to (7, 3)} \\
            }
        }
    }
    & & \\ \\

    \theutterance \stepcounter{utterance}  
    & & \multicolumn{4}{p{0.6\linewidth}}{
        \cellcolor[rgb]{0.9,0.9,0.9}{
            \makecell[{{p{\linewidth}}}]{
                \texttt{\tiny{[GM$\rangle$P2]}}
                \texttt{تم تمرير رسالتك إلى اللاعب الآخر.} \\
\\ 
\\ 
\texttt{أنت تلعب حالياً الجولة 16 من أصل 20 كحد أقصى.} \\
\texttt{لقد جمعتما معاً 10 من أصل 12 عقوبة.} \\
\texttt{حرّك اللاعب الآخر كائناً على شبكته.} \\
\\ 
\\ 
\texttt{ما هو أمرك التالي؟} \\
            }
        }
    }
    & \\ \\

    \theutterance \stepcounter{utterance}  
    & & & \multicolumn{4}{p{0.6\linewidth}}{
        \cellcolor[rgb]{0.9,0.9,0.9}{
            \makecell[{{p{\linewidth}}}]{
                \texttt{\tiny{[GM$\langle$P2]}}
                \texttt{قل: سأراقب تحريك اللاعب الآخر، ثم سأحدد الخطوة المقابلة على لوحتي لتحقيق هدف المحاذاة.} \\
            }
        }
    }
     \\ \\

    \theutterance \stepcounter{utterance}  
    & & & \multicolumn{2}{p{0.3\linewidth}}{
        \cellcolor[rgb]{0.9,0.9,0.9}{
            \makecell[{{p{\linewidth}}}]{
                \texttt{\tiny{[GM$|$GM]}}
                \texttt{قل: سأراقب تحريك اللاعب الآخر، ثم سأحدد الخطوة المقابلة على لوحتي لتحقيق هدف المحاذاة.} \\
            }
        }
    }
    & & \\ \\

    \theutterance \stepcounter{utterance}  
    & & \multicolumn{4}{p{0.6\linewidth}}{
        \cellcolor[rgb]{0.9,0.9,0.9}{
            \makecell[{{p{\linewidth}}}]{
                \texttt{\tiny{[P1$\langle$GM]}}
                \texttt{تم تحريك 'W' إلى (7,3) بنجاح. تبدو شبكتك المحدّثة هكذا:} \\
\texttt{```} \\
\texttt{ 1234567} \\
\texttt{╔══╤════╗} \\
\texttt{║◌◌│◌◌◌H║ 1} \\
\texttt{╟──┤◌◌┌─╢ 2} \\
\texttt{║CI├──┤W║ 3} \\
\texttt{╟──┘◌◌├─╢ 4} \\
\texttt{║◌◌◌◌◌│◌║ 5} \\
\texttt{╟─┐◌┌─┼─╢ 6} \\
\texttt{║◌│T│◌│◌║ 7} \\
\texttt{╚═╧═╧═╧═╝} \\
\texttt{```} \\
\texttt{أنت تلعب حالياً الجولة 16 من أصل 20 كحد أقصى.} \\
\texttt{لقد جمعتما معاً 10 من أصل 12 عقوبة.} \\
\texttt{أرسل اللاعب الآخر هذه الرسالة:} \\
\\ 
\\ 
\texttt{"سأراقب تحريك اللاعب الآخر، ثم سأحدد الخطوة المقابلة على لوحتي لتحقيق هدف المحاذاة."} \\
\\ 
\\ 
\texttt{ما هو أمرك التالي؟} \\
            }
        }
    }
    & \\ \\

    \theutterance \stepcounter{utterance}  
    & \multicolumn{4}{p{0.6\linewidth}}{
        \cellcolor[rgb]{0.9,0.9,0.9}{
            \makecell[{{p{\linewidth}}}]{
                \texttt{\tiny{[P1$\rangle$GM]}}
                \texttt{قل: شكراً لك على الاستمرارية. سأراقب التحريك الذي قام به اللاعب الآخر، ثم سأحدد الخطوة المقابلة على لوحتي لتحقيق الهدف المشترك للمحاذاة. يرجى التأكيد أنني فهمت بشكل صحيح؟} \\
            }
        }
    }
    & & \\ \\

    \theutterance \stepcounter{utterance}  
    & & & \multicolumn{2}{p{0.3\linewidth}}{
        \cellcolor[rgb]{0.9,0.9,0.9}{
            \makecell[{{p{\linewidth}}}]{
                \texttt{\tiny{[GM$|$GM]}}
                \texttt{قل: شكراً لك على الاستمرارية. سأراقب التحريك الذي قام به اللاعب الآخر، ثم سأحدد الخطوة المقابلة على لوحتي لتحقيق الهدف المشترك للمحاذاة. يرجى التأكيد أنني فهمت بشكل صحيح؟} \\
            }
        }
    }
    & & \\ \\

    \theutterance \stepcounter{utterance}  
    & & \multicolumn{4}{p{0.6\linewidth}}{
        \cellcolor[rgb]{0.9,0.9,0.9}{
            \makecell[{{p{\linewidth}}}]{
                \texttt{\tiny{[GM$\rangle$P2]}}
                \texttt{تم تمرير رسالتك إلى اللاعب الآخر.} \\
\\ 
\\ 
\texttt{أنت تلعب حالياً الجولة 17 من أصل 20 كحد أقصى.} \\
\texttt{لقد جمعتما معاً 10 من أصل 12 عقوبة.} \\
\texttt{أرسل اللاعب الآخر هذه الرسالة:} \\
\\ 
\\ 
\texttt{"شكراً لك على الاستمرارية. سأراقب التحريك الذي قام به اللاعب الآخر، ثم سأحدد الخطوة المقابلة على لوحتي لتحقيق الهدف المشترك للمحاذاة. يرجى التأكيد أنني فهمت بشكل صحيح؟"} \\
\\ 
\\ 
\texttt{ما هو أمرك التالي؟} \\
            }
        }
    }
    & \\ \\

    \theutterance \stepcounter{utterance}  
    & & & \multicolumn{4}{p{0.6\linewidth}}{
        \cellcolor[rgb]{0.9,0.9,0.9}{
            \makecell[{{p{\linewidth}}}]{
                \texttt{\tiny{[GM$\langle$P2]}}
                \texttt{قل: أنا أؤكد لك فهمك الصحيح للخطوات المقابلة: إذا تحرك كائن معين على لوحته إلى إحداثيات معينة، فأنا سأتحرك الكائن المقابل على لوحتي إلى نفس الإحداثيات. إذا كان هذا صحيحاً، أنا مستعد لتنفيذ حركتي الموافقة الآن.} \\
            }
        }
    }
     \\ \\

    \theutterance \stepcounter{utterance}  
    & & & \multicolumn{2}{p{0.3\linewidth}}{
        \cellcolor[rgb]{0.9,0.9,0.9}{
            \makecell[{{p{\linewidth}}}]{
                \texttt{\tiny{[GM$|$GM]}}
                \texttt{قل: أنا أؤكد لك فهمك الصحيح للخطوات المقابلة: إذا تحرك كائن معين على لوحته إلى إحداثيات معينة، فأنا سأتحرك الكائن المقابل على لوحتي إلى نفس الإحداثيات. إذا كان هذا صحيحاً، أنا مستعد لتنفيذ حركتي الموافقة الآن.} \\
            }
        }
    }
    & & \\ \\

    \theutterance \stepcounter{utterance}  
    & & \multicolumn{4}{p{0.6\linewidth}}{
        \cellcolor[rgb]{0.9,0.9,0.9}{
            \makecell[{{p{\linewidth}}}]{
                \texttt{\tiny{[P1$\langle$GM]}}
                \texttt{تم تمرير رسالتك إلى اللاعب الآخر.} \\
\\ 
\\ 
\texttt{أنت تلعب حالياً الجولة 17 من أصل 20 كحد أقصى.} \\
\texttt{لقد جمعتما معاً 10 من أصل 12 عقوبة.} \\
\texttt{أرسل اللاعب الآخر هذه الرسالة:} \\
\\ 
\\ 
\texttt{"أنا أؤكد لك فهمك الصحيح للخطوات المقابلة: إذا تحرك كائن معين على لوحته إلى إحداثيات معينة، فأنا سأتحرك الكائن المقابل على لوحتي إلى نفس الإحداثيات. إذا كان هذا صحيحاً، أنا مستعد لتنفيذ حركتي الموافقة الآن."} \\
\\ 
\\ 
\texttt{ما هو أمرك التالي؟} \\
            }
        }
    }
    & \\ \\

    \theutterance \stepcounter{utterance}  
    & \multicolumn{4}{p{0.6\linewidth}}{
        \cellcolor[rgb]{0.9,0.9,0.9}{
            \makecell[{{p{\linewidth}}}]{
                \texttt{\tiny{[P1$\rangle$GM]}}
                \texttt{حرّك: H, (2,2)} \\
            }
        }
    }
    & & \\ \\

    \theutterance \stepcounter{utterance}  
    & & & \multicolumn{2}{p{0.3\linewidth}}{
        \cellcolor[rgb]{0.9,0.9,0.9}{
            \makecell[{{p{\linewidth}}}]{
                \texttt{\tiny{[GM$|$GM]}}
                \texttt{حرّك: H, (2,2)} \\
            }
        }
    }
    & & \\ \\

    \theutterance \stepcounter{utterance}  
    & & & \multicolumn{2}{p{0.3\linewidth}}{
        \cellcolor[rgb]{0.9,0.9,0.9}{
            \makecell[{{p{\linewidth}}}]{
                \texttt{\tiny{[GM$|$GM]}}
                \texttt{عقوبة: (2,2) ليست فارغة، بل تحتوي على '─'.} \\
\\ 
\texttt{لقد جمعتما معاً 11 من أصل 12 عقوبة.} \\
\texttt{يرجى المحاولة مرة أخرى!} \\
            }
        }
    }
    & & \\ \\

    \theutterance \stepcounter{utterance}  
    & & \multicolumn{4}{p{0.6\linewidth}}{
        \cellcolor[rgb]{0.9,0.9,0.9}{
            \makecell[{{p{\linewidth}}}]{
                \texttt{\tiny{[P1$\langle$GM]}}
                \texttt{عقوبة: (2,2) ليست فارغة، بل تحتوي على '─'.} \\
\\ 
\texttt{لقد جمعتما معاً 11 من أصل 12 عقوبة.} \\
\texttt{يرجى المحاولة مرة أخرى!} \\
            }
        }
    }
    & \\ \\

    \theutterance \stepcounter{utterance}  
    & \multicolumn{4}{p{0.6\linewidth}}{
        \cellcolor[rgb]{0.9,0.9,0.9}{
            \makecell[{{p{\linewidth}}}]{
                \texttt{\tiny{[P1$\rangle$GM]}}
                \texttt{حرّك: H, (2,4)} \\
            }
        }
    }
    & & \\ \\

    \theutterance \stepcounter{utterance}  
    & & & \multicolumn{2}{p{0.3\linewidth}}{
        \cellcolor[rgb]{0.9,0.9,0.9}{
            \makecell[{{p{\linewidth}}}]{
                \texttt{\tiny{[GM$|$GM]}}
                \texttt{حرّك: H, (2,4)} \\
            }
        }
    }
    & & \\ \\

    \theutterance \stepcounter{utterance}  
    & & & \multicolumn{2}{p{0.3\linewidth}}{
        \cellcolor[rgb]{0.9,0.9,0.9}{
            \makecell[{{p{\linewidth}}}]{
                \texttt{\tiny{[GM$|$GM]}}
                \texttt{عقوبة: (2,4) ليست فارغة، بل تحتوي على '─'.} \\
\\ 
\texttt{لقد جمعتما معاً 12 من أصل 12 عقوبة.} \\
\texttt{يرجى المحاولة مرة أخرى!} \\
            }
        }
    }
    & & \\ \\

    \theutterance \stepcounter{utterance}  
    & & \multicolumn{4}{p{0.6\linewidth}}{
        \cellcolor[rgb]{0.9,0.9,0.9}{
            \makecell[{{p{\linewidth}}}]{
                \texttt{\tiny{[P1$\langle$GM]}}
                \texttt{عقوبة: (2,4) ليست فارغة، بل تحتوي على '─'.} \\
\\ 
\texttt{لقد جمعتما معاً 12 من أصل 12 عقوبة.} \\
\texttt{يرجى المحاولة مرة أخرى!} \\
            }
        }
    }
    & \\ \\

    \theutterance \stepcounter{utterance}  
    & \multicolumn{4}{p{0.6\linewidth}}{
        \cellcolor[rgb]{0.9,0.9,0.9}{
            \makecell[{{p{\linewidth}}}]{
                \texttt{\tiny{[P1$\rangle$GM]}}
                \texttt{قل: لقد وصلنا إلى 12 عقوبة، أي الحد الأقصى. سنقوم بإعادة النظر في الخطة ونحدد إستراتيجية جديدة. هل يمكننا إعادة جدولة التحريك بطريقة مختلفة؟} \\
            }
        }
    }
    & & \\ \\

    \theutterance \stepcounter{utterance}  
    & & & \multicolumn{2}{p{0.3\linewidth}}{
        \cellcolor[rgb]{0.9,0.9,0.9}{
            \makecell[{{p{\linewidth}}}]{
                \texttt{\tiny{[GM$|$GM]}}
                \texttt{قل: لقد وصلنا إلى 12 عقوبة، أي الحد الأقصى. سنقوم بإعادة النظر في الخطة ونحدد إستراتيجية جديدة. هل يمكننا إعادة جدولة التحريك بطريقة مختلفة؟} \\
            }
        }
    }
    & & \\ \\

    \theutterance \stepcounter{utterance}  
    & & \multicolumn{4}{p{0.6\linewidth}}{
        \cellcolor[rgb]{0.9,0.9,0.9}{
            \makecell[{{p{\linewidth}}}]{
                \texttt{\tiny{[GM$\rangle$P2]}}
                \texttt{تم تمرير رسالتك إلى اللاعب الآخر.} \\
\\ 
\\ 
\texttt{أنت تلعب حالياً الجولة 18 من أصل 20 كحد أقصى.} \\
\texttt{لقد جمعتما معاً 12 من أصل 12 عقوبة.} \\
\texttt{أرسل اللاعب الآخر هذه الرسالة:} \\
\\ 
\\ 
\texttt{"لقد وصلنا إلى 12 عقوبة، أي الحد الأقصى. سنقوم بإعادة النظر في الخطة ونحدد إستراتيجية جديدة. هل يمكننا إعادة جدولة التحريك بطريقة مختلفة؟"} \\
\\ 
\\ 
\texttt{ما هو أمرك التالي؟} \\
            }
        }
    }
    & \\ \\

    \theutterance \stepcounter{utterance}  
    & & & \multicolumn{4}{p{0.6\linewidth}}{
        \cellcolor[rgb]{0.9,0.9,0.9}{
            \makecell[{{p{\linewidth}}}]{
                \texttt{\tiny{[GM$\langle$P2]}}
                \texttt{قل: نتفق على إعادة جدولة التحريك بطريقة مختلفة لتجنب تجاوز الحد الأقصى للعقوبات. سأقدم اقتراح استراتيجي جديد لبناء الخطة المشتركة، وسأطلب رأيك وتعاونك لتنفيذه.} \\
            }
        }
    }
     \\ \\

    \theutterance \stepcounter{utterance}  
    & & & \multicolumn{2}{p{0.3\linewidth}}{
        \cellcolor[rgb]{0.9,0.9,0.9}{
            \makecell[{{p{\linewidth}}}]{
                \texttt{\tiny{[GM$|$GM]}}
                \texttt{قل: نتفق على إعادة جدولة التحريك بطريقة مختلفة لتجنب تجاوز الحد الأقصى للعقوبات. سأقدم اقتراح استراتيجي جديد لبناء الخطة المشتركة، وسأطلب رأيك وتعاونك لتنفيذه.} \\
            }
        }
    }
    & & \\ \\

    \theutterance \stepcounter{utterance}  
    & & \multicolumn{4}{p{0.6\linewidth}}{
        \cellcolor[rgb]{0.9,0.9,0.9}{
            \makecell[{{p{\linewidth}}}]{
                \texttt{\tiny{[P1$\langle$GM]}}
                \texttt{تم تمرير رسالتك إلى اللاعب الآخر.} \\
\\ 
\\ 
\texttt{أنت تلعب حالياً الجولة 18 من أصل 20 كحد أقصى.} \\
\texttt{لقد جمعتما معاً 12 من أصل 12 عقوبة.} \\
\texttt{أرسل اللاعب الآخر هذه الرسالة:} \\
\\ 
\\ 
\texttt{"نتفق على إعادة جدولة التحريك بطريقة مختلفة لتجنب تجاوز الحد الأقصى للعقوبات. سأقدم اقتراح استراتيجي جديد لبناء الخطة المشتركة، وسأطلب رأيك وتعاونك لتنفيذه."} \\
\\ 
\\ 
\texttt{ما هو أمرك التالي؟} \\
            }
        }
    }
    & \\ \\

    \theutterance \stepcounter{utterance}  
    & \multicolumn{4}{p{0.6\linewidth}}{
        \cellcolor[rgb]{0.9,0.9,0.9}{
            \makecell[{{p{\linewidth}}}]{
                \texttt{\tiny{[P1$\rangle$GM]}}
                \texttt{قل: شكراً لك على التصريح بإعادة جدولة التحريك. سأقوم بوضع خطة استراتيجية جديدة تعتمد على تجنّب العوائق التي واجهناها، وأعدّ خطوات التحريك لكل كائن بطريقة مختلفة. يرجى الانتظار لباقي التحقيقات في الاستراتيجية الجديدة.} \\
            }
        }
    }
    & & \\ \\

    \theutterance \stepcounter{utterance}  
    & & & \multicolumn{2}{p{0.3\linewidth}}{
        \cellcolor[rgb]{0.9,0.9,0.9}{
            \makecell[{{p{\linewidth}}}]{
                \texttt{\tiny{[GM$|$GM]}}
                \texttt{قل: شكراً لك على التصريح بإعادة جدولة التحريك. سأقوم بوضع خطة استراتيجية جديدة تعتمد على تجنّب العوائق التي واجهناها، وأعدّ خطوات التحريك لكل كائن بطريقة مختلفة. يرجى الانتظار لباقي التحقيقات في الاستراتيجية الجديدة.} \\
            }
        }
    }
    & & \\ \\

    \theutterance \stepcounter{utterance}  
    & & \multicolumn{4}{p{0.6\linewidth}}{
        \cellcolor[rgb]{0.9,0.9,0.9}{
            \makecell[{{p{\linewidth}}}]{
                \texttt{\tiny{[GM$\rangle$P2]}}
                \texttt{تم تمرير رسالتك إلى اللاعب الآخر.} \\
\\ 
\\ 
\texttt{أنت تلعب حالياً الجولة 19 من أصل 20 كحد أقصى.} \\
\texttt{لقد جمعتما معاً 12 من أصل 12 عقوبة.} \\
\texttt{أرسل اللاعب الآخر هذه الرسالة:} \\
\\ 
\\ 
\texttt{"شكراً لك على التصريح بإعادة جدولة التحريك. سأقوم بوضع خطة استراتيجية جديدة تعتمد على تجنّب العوائق التي واجهناها، وأعدّ خطوات التحريك لكل كائن بطريقة مختلفة. يرجى الانتظار لباقي التحقيقات في الاستراتيجية الجديدة."} \\
\\ 
\\ 
\texttt{ما هو أمرك التالي؟} \\
            }
        }
    }
    & \\ \\

    \theutterance \stepcounter{utterance}  
    & & & \multicolumn{4}{p{0.6\linewidth}}{
        \cellcolor[rgb]{0.9,0.9,0.9}{
            \makecell[{{p{\linewidth}}}]{
                \texttt{\tiny{[GM$\langle$P2]}}
                \texttt{قل: أنا متقبل للخطة الاستراتيجية الجديدة. سأراقب تنفيذ الخطوات التي تقترحها لإعادة جدولة التحريك وتجنب العوائق، وسأشارك في تحليل الخطوات المقابلة على لوحتي لتحقيق الهدف المشترك للمحاذاة.} \\
            }
        }
    }
     \\ \\

    \theutterance \stepcounter{utterance}  
    & & & \multicolumn{2}{p{0.3\linewidth}}{
        \cellcolor[rgb]{0.9,0.9,0.9}{
            \makecell[{{p{\linewidth}}}]{
                \texttt{\tiny{[GM$|$GM]}}
                \texttt{قل: أنا متقبل للخطة الاستراتيجية الجديدة. سأراقب تنفيذ الخطوات التي تقترحها لإعادة جدولة التحريك وتجنب العوائق، وسأشارك في تحليل الخطوات المقابلة على لوحتي لتحقيق الهدف المشترك للمحاذاة.} \\
            }
        }
    }
    & & \\ \\

    \theutterance \stepcounter{utterance}  
    & & \multicolumn{4}{p{0.6\linewidth}}{
        \cellcolor[rgb]{0.9,0.9,0.9}{
            \makecell[{{p{\linewidth}}}]{
                \texttt{\tiny{[P1$\langle$GM]}}
                \texttt{تم تمرير رسالتك إلى اللاعب الآخر.} \\
\\ 
\\ 
\texttt{أنت تلعب حالياً الجولة 19 من أصل 20 كحد أقصى.} \\
\texttt{لقد جمعتما معاً 12 من أصل 12 عقوبة.} \\
\texttt{أرسل اللاعب الآخر هذه الرسالة:} \\
\\ 
\\ 
\texttt{"أنا متقبل للخطة الاستراتيجية الجديدة. سأراقب تنفيذ الخطوات التي تقترحها لإعادة جدولة التحريك وتجنب العوائق، وسأشارك في تحليل الخطوات المقابلة على لوحتي لتحقيق الهدف المشترك للمحاذاة."} \\
\\ 
\\ 
\texttt{ما هو أمرك التالي؟} \\
            }
        }
    }
    & \\ \\

    \theutterance \stepcounter{utterance}  
    & \multicolumn{4}{p{0.6\linewidth}}{
        \cellcolor[rgb]{0.9,0.9,0.9}{
            \makecell[{{p{\linewidth}}}]{
                \texttt{\tiny{[P1$\rangle$GM]}}
                \texttt{قل: شكراً لك على القبول بالخطة الاستراتيجية الجديدة. سأقوم بوضع خطوات التحريك لكل كائن بطريقة مختلفة لتجنب العوائق، وسأحدد الخطوات المقابلة على لوحتي. يرجى التأكيد أنني فهمت خطوات التحريك الصحيحة التي تخطط لتنفيذها؟} \\
            }
        }
    }
    & & \\ \\

    \theutterance \stepcounter{utterance}  
    & & & \multicolumn{2}{p{0.3\linewidth}}{
        \cellcolor[rgb]{0.9,0.9,0.9}{
            \makecell[{{p{\linewidth}}}]{
                \texttt{\tiny{[GM$|$GM]}}
                \texttt{قل: شكراً لك على القبول بالخطة الاستراتيجية الجديدة. سأقوم بوضع خطوات التحريك لكل كائن بطريقة مختلفة لتجنب العوائق، وسأحدد الخطوات المقابلة على لوحتي. يرجى التأكيد أنني فهمت خطوات التحريك الصحيحة التي تخطط لتنفيذها؟} \\
            }
        }
    }
    & & \\ \\

    \theutterance \stepcounter{utterance}  
    & & \multicolumn{4}{p{0.6\linewidth}}{
        \cellcolor[rgb]{0.9,0.9,0.9}{
            \makecell[{{p{\linewidth}}}]{
                \texttt{\tiny{[GM$\rangle$P2]}}
                \texttt{تم تمرير رسالتك إلى اللاعب الآخر.} \\
\\ 
\\ 
\texttt{أنت تلعب حالياً الجولة 20 من أصل 20 كحد أقصى.} \\
\texttt{لقد جمعتما معاً 12 من أصل 12 عقوبة.} \\
\texttt{أرسل اللاعب الآخر هذه الرسالة:} \\
\\ 
\\ 
\texttt{"شكراً لك على القبول بالخطة الاستراتيجية الجديدة. سأقوم بوضع خطوات التحريك لكل كائن بطريقة مختلفة لتجنب العوائق، وسأحدد الخطوات المقابلة على لوحتي. يرجى التأكيد أنني فهمت خطوات التحريك الصحيحة التي تخطط لتنفيذها؟"} \\
\\ 
\\ 
\texttt{ما هو أمرك التالي؟} \\
            }
        }
    }
    & \\ \\

    \theutterance \stepcounter{utterance}  
    & & & \multicolumn{4}{p{0.6\linewidth}}{
        \cellcolor[rgb]{0.9,0.9,0.9}{
            \makecell[{{p{\linewidth}}}]{
                \texttt{\tiny{[GM$\langle$P2]}}
                \texttt{قل: أنا أؤكد لكم فهمي لكافة خطوات التحريك التي حررتموها لتجنب العوائق وإكمال المهمة. إذا كانت هناك أي تفاصيل تحتاج لتوضيح أو تأكيد دقيق، فأنا مستعد لأكرر التأكيد أو إضافة أي توضيحات ضرورية لضمان وضوح الخطة النهائية.} \\
            }
        }
    }
     \\ \\

    \theutterance \stepcounter{utterance}  
    & & & \multicolumn{2}{p{0.3\linewidth}}{
        \cellcolor[rgb]{0.9,0.9,0.9}{
            \makecell[{{p{\linewidth}}}]{
                \texttt{\tiny{[GM$|$GM]}}
                \texttt{قل: أنا أؤكد لكم فهمي لكافة خطوات التحريك التي حررتموها لتجنب العوائق وإكمال المهمة. إذا كانت هناك أي تفاصيل تحتاج لتوضيح أو تأكيد دقيق، فأنا مستعد لأكرر التأكيد أو إضافة أي توضيحات ضرورية لضمان وضوح الخطة النهائية.} \\
            }
        }
    }
    & & \\ \\

    \theutterance \stepcounter{utterance}  
    & & & \multicolumn{2}{p{0.3\linewidth}}{
        \cellcolor[rgb]{0.9,0.9,0.9}{
            \makecell[{{p{\linewidth}}}]{
                \texttt{\tiny{[GM$|$GM]}}
                \texttt{Maximum number of rounds reached} \\
            }
        }
    }
    & & \\ \\

    \theutterance \stepcounter{utterance}  
    & & & \multicolumn{2}{p{0.3\linewidth}}{
        \cellcolor[rgb]{0.9,0.9,0.9}{
            \makecell[{{p{\linewidth}}}]{
                \texttt{\tiny{[GM$|$GM]}}
                \texttt{Initial states:} \\
\texttt{```} \\
\texttt{Player 1:} \\
\texttt{ 1234567} \\
\texttt{╔══╤════╗} \\
\texttt{║◌◌│◌◌◌H║ 1} \\
\texttt{╟──┤◌◌┌─╢ 2} \\
\texttt{║CI├──┤◌║ 3} \\
\texttt{╟──┘◌◌├─╢ 4} \\
\texttt{║◌◌◌◌◌│◌║ 5} \\
\texttt{╟─┐W┌─┼─╢ 6} \\
\texttt{║◌│T│◌│◌║ 7} \\
\texttt{╚═╧═╧═╧═╝} \\
\\ 
\\ 
\texttt{Player 2:} \\
\texttt{ 1234567} \\
\texttt{╔══╤════╗} \\
\texttt{║◌W│◌◌◌◌║ 1} \\
\texttt{╟──┤◌H┌─╢ 2} \\
\texttt{║◌◌├──┤◌║ 3} \\
\texttt{╟──┘◌◌├─╢ 4} \\
\texttt{║◌◌◌◌◌│◌║ 5} \\
\texttt{╟─┐C┌─┼─╢ 6} \\
\texttt{║I│◌│T│◌║ 7} \\
\texttt{╚═╧═╧═╧═╝} \\
\\ 
\texttt{```} \\
            }
        }
    }
    & & \\ \\

    \theutterance \stepcounter{utterance}  
    & & & \multicolumn{2}{p{0.3\linewidth}}{
        \cellcolor[rgb]{0.9,0.9,0.9}{
            \makecell[{{p{\linewidth}}}]{
                \texttt{\tiny{[GM$|$GM]}}
                \texttt{End states:} \\
\texttt{```} \\
\texttt{Player 1:} \\
\texttt{ 1234567} \\
\texttt{╔══╤════╗} \\
\texttt{║◌◌│◌◌◌H║ 1} \\
\texttt{╟──┤◌◌┌─╢ 2} \\
\texttt{║CI├──┤W║ 3} \\
\texttt{╟──┘◌◌├─╢ 4} \\
\texttt{║◌◌◌◌◌│◌║ 5} \\
\texttt{╟─┐◌┌─┼─╢ 6} \\
\texttt{║◌│T│◌│◌║ 7} \\
\texttt{╚═╧═╧═╧═╝} \\
\\ 
\\ 
\texttt{Player 2:} \\
\texttt{ 1234567} \\
\texttt{╔══╤════╗} \\
\texttt{║◌W│◌◌◌◌║ 1} \\
\texttt{╟──┤◌H┌─╢ 2} \\
\texttt{║◌◌├──┤◌║ 3} \\
\texttt{╟──┘◌◌├─╢ 4} \\
\texttt{║◌◌◌◌◌│◌║ 5} \\
\texttt{╟─┐C┌─┼─╢ 6} \\
\texttt{║I│◌│T│◌║ 7} \\
\texttt{╚═╧═╧═╧═╝} \\
\\ 
\texttt{```} \\
            }
        }
    }
    & & \\ \\

    \theutterance \stepcounter{utterance}  
    & & & \multicolumn{2}{p{0.3\linewidth}}{
        \cellcolor[rgb]{0.9,0.9,0.9}{
            \makecell[{{p{\linewidth}}}]{
                \texttt{\tiny{[GM$|$GM]}}
                \texttt{* success: True} \\
\texttt{* lose: False} \\
\texttt{* aborted: False} \\
\texttt{{-}{-}{-}{-}{-}{-}{-}} \\
\texttt{* Shifts: 0.00} \\
\texttt{* Max Shifts: 8.00} \\
\texttt{* Min Shifts: 4.00} \\
\texttt{* End Distance Sum: 17.35} \\
\texttt{* Init Distance Sum: 17.06} \\
\texttt{* Expected Distance Sum: 20.95} \\
\texttt{* Penalties: 12.00} \\
\texttt{* Max Penalties: 12.00} \\
\texttt{* Rounds: 20.00} \\
\texttt{* Max Rounds: 20.00} \\
\texttt{* Object Count: 5.00} \\
            }
        }
    }
    & & \\ \\

\end{supertabular}
}

\end{document}
